% explainer-source-hash: b61f9ba78915fdf5b2011a630b8496641f14462f328b2f9ba5f40a7d1eb57aaa
% explainer-source-commit: 4302e8d356b88b447a7e8b3cbbfab75e1d7e0dc8
% explainer-generated: 2026-07-16
% Regenerate with /explain-all (see WIRING.md). Do not edit this stamp by hand:
% it is how the toolchain knows whether this document still matches the code.
\documentclass[11pt,a4paper]{article}
\input{preamble}

\title{\vspace{-1.5cm}\textbf{raft-kv}\\[0.3cm]
\large A Distributed Key-Value Store with Raft Consensus,\\
Deterministic Simulation, and a Linearizability Checker\\[0.4cm]
\normalsize Deep Dive: an exhaustive, first-principles walkthrough}
\author{}
\date{}

\begin{document}
\maketitle
\thispagestyle{empty}

\begin{keyidea}{Read this first: what the project is, in one paragraph}
\code{raft-kv} is a database that stores simple key-value pairs (like
\code{city = lahore}) and keeps identical copies of that data on several
computers at once, so that the data survives any one of them dying. Keeping
several computers in exact agreement, while they crash and their network
breaks, is the problem called \emph{consensus}. This project solves it with an
algorithm called \emph{Raft}, written from scratch in Python with no consensus
libraries. But the code that solves consensus is only about a third of the
project. The other two thirds are the interesting part: a \emph{simulator} that
can replay thousands of randomized disasters against the real consensus code
in a way that is exactly reproducible from a single integer, and a
\emph{linearizability checker} that mathematically verifies that the answers
clients received could have come from a single, honest, non-distributed
database.
\end{keyidea}

\tableofcontents
\newpage

% =====================================================================
\section{Orientation: the problem, from zero}

\subsection{Why one computer is not enough}

Suppose you store data on one machine. A single hard disk failure, a power cut,
or a careless \code{kill} command and the data is gone or unreachable. The
obvious fix is to keep copies on several machines. That word, copies, is where
all the difficulty starts.

\begin{jargon}{Replication}
Keeping more than one copy of the same data on more than one machine, so that
the loss of any single machine does not lose the data. Each copy is called a
\emph{replica}. The machines are also called \emph{nodes} or \emph{servers}.
\end{jargon}

\begin{jargon}{Node}
One participating computer (or, in this project's tests, one process, or even
one Python object). In \code{raft-kv} nodes have short string names:
\code{n1}, \code{n2}, \code{n3}.
\end{jargon}

The naive design is: whenever a write arrives, send it to all three machines.
That fails immediately, for reasons worth stating concretely:

\begin{itemize}
\item Machine \code{n3} is briefly down when the write arrives, so it misses
      it. Its copy is now wrong, and nothing in the system knows.
\item Two clients write different values to the same key at the same moment.
      Machine \code{n1} happens to process them in one order, \code{n2} in the
      other. The two copies now permanently disagree.
\item The network splits in half. Machines on each side think the other side
      died, and both sides keep accepting writes. Now there are two divergent
      databases claiming to be the same database. This is called
      \emph{split brain}, and it is the disease that consensus exists to
      prevent.
\end{itemize}

\begin{jargon}{Split brain}
Two or more sets of nodes each believe they are the authoritative cluster, and
each accepts writes independently. When the network heals, there is no honest
way to merge the two divergent versions of history: acknowledged writes must be
thrown away. Any correct consensus algorithm makes split brain structurally
impossible rather than merely unlikely.
\end{jargon}

\subsection{Consensus, from first principles}

\begin{keyidea}{What consensus actually means}
Consensus is the problem of getting a group of unreliable machines to agree on
one \textbf{ordered list of decisions}, such that once a decision is agreed, it
is agreed forever, by everyone, and no two machines ever disagree about what
decision number 7 was. That is the whole idea. Everything else in Raft is
machinery to make that safe under crashes and network failure.
\end{keyidea}

Notice the pivot: the group does not agree on ``the value of the key
\code{city}''. It agrees on an ordered \emph{list of commands}:

\begin{center}
\begin{tabular}{cl}
\toprule
\textbf{Position} & \textbf{Command} \\
\midrule
1 & \code{put city = lahore} \\
2 & \code{put team = raptors} \\
3 & \code{cas team: raptors $\rightarrow$ dinos} \\
4 & \code{delete city} \\
\bottomrule
\end{tabular}
\end{center}

If every machine agrees on that list, and every machine executes the list in
order starting from an empty database, then every machine ends up with an
identical database, necessarily. Agreement on the data is derived for free from
agreement on the list. This trick has a name.

\begin{jargon}{Replicated state machine}
A design where each node holds (a) an identical ordered log of commands and
(b) a deterministic program that executes those commands in order. Because the
program is deterministic (same inputs, same order, same outputs, always) and
the log is identical, the resulting state is identical on every node. The hard
problem is reduced to: keep the logs identical. In this project the log is
\code{raft/log.py}, and the deterministic program is \code{kv/machine.py}.
\end{jargon}

\begin{figure}[H]
\centering
\resizebox{\textwidth}{!}{
\begin{tikzpicture}[node distance=6mm]
\node[box, minimum width=30mm] (c) {Client\\\scriptsize{}put city = lahore};

\node[draw=inkblue, fill=softgrey, rounded corners=3pt,
      minimum width=42mm, minimum height=34mm, right=22mm of c] (n1f) {};
\node[lbl, anchor=north] at (n1f.north) {node n1 (leader)};
\node[box, minimum width=34mm] (l1) at ($(n1f.center)+(0,3mm)$)
      {log: [1:put city, 2:...]};
\node[store, minimum width=24mm] (m1) at ($(n1f.center)+(0,-10mm)$)
      {\scriptsize state machine};

\node[draw=linegrey, fill=white, rounded corners=3pt,
      minimum width=42mm, minimum height=34mm, right=16mm of n1f] (n2f) {};
\node[lbl, anchor=north] at (n2f.north) {node n2 (follower)};
\node[box, minimum width=34mm] (l2) at ($(n2f.center)+(0,3mm)$)
      {log: [1:put city, 2:...]};
\node[store, minimum width=24mm] (m2) at ($(n2f.center)+(0,-10mm)$)
      {\scriptsize state machine};

\node[draw=linegrey, fill=white, rounded corners=3pt,
      minimum width=42mm, minimum height=34mm, right=16mm of n2f] (n3f) {};
\node[lbl, anchor=north] at (n3f.north) {node n3 (follower)};
\node[box, minimum width=34mm] (l3) at ($(n3f.center)+(0,3mm)$)
      {log: [1:put city, 2:...]};
\node[store, minimum width=24mm] (m3) at ($(n3f.center)+(0,-10mm)$)
      {\scriptsize state machine};

\draw[flow] (c) -- (n1f);
\draw[flow] (l1) -- node[lbl, above]{replicate} (l2);
\draw[flow] (l2) -- node[lbl, above]{replicate} (l3);
\draw[flow] (l1) -- node[lbl, right]{apply in order} (m1);
\draw[flow] (l2) -- node[lbl, right]{apply in order} (m2);
\draw[flow] (l3) -- node[lbl, right]{apply in order} (m3);
\end{tikzpicture}}
\caption{The replicated state machine. Agreement on the ordered log makes
agreement on the data automatic, because applying the same commands in the same
order to the same starting state must produce the same result.}
\end{figure}

\subsection{The majority rule: why three, not two}

The single mechanism that makes split brain impossible is deceptively small.

\begin{keyidea}{Quorums, and why they work}
Every important decision requires the agreement of a strict majority of the
cluster (a \emph{quorum}): 2 out of 3, 3 out of 5. The reason this works is a
fact about sets, not about computers: \textbf{any two majorities of the same
set must share at least one member}. Two majorities of \{n1, n2, n3\} cannot be
disjoint, because 2 + 2 > 3. So if a decision needs a majority, and some node
in every majority remembers the previous decision, no second, contradictory
decision can slip through unseen. A partition can therefore leave at most
\emph{one} side with a majority; the other side is mathematically unable to act.
\end{keyidea}

\begin{jargon}{Quorum / majority}
More than half the nodes in the current cluster configuration. In this code it
is one line, \code{raft/node.py}: \code{len(self.config) // 2 + 1}. For 3 nodes
that is 2; for 4 nodes it is 3, which is why even-sized clusters buy you
nothing: a 4-node cluster tolerates exactly one failure, the same as a 3-node
cluster, while being more expensive and slower.
\end{jargon}

This is why the project's self-test (section~\ref{sec:canary}) attacks
precisely this line. Weaken \code{// 2 + 1} into \code{// 2} and the majority
overlap guarantee evaporates, split brain becomes possible, and the whole
edifice falls. It is the load-bearing wall.

\subsection{Linearizability: the correctness target}

Replication is not the goal. Replication that behaves \emph{as if it were not
replicated} is the goal.

\begin{jargon}{Linearizability}
The strongest single-object consistency guarantee. A system is linearizable if
every operation appears to take effect \textbf{instantaneously at some single
moment between the client sending the request and receiving the response}, and
that ordering of moments is consistent with real time. The practical meaning:
if your write finished at 10:00:01, then \emph{every} read that starts at
10:00:02, on any node, from any client, must see it. The distributed system is
indistinguishable, from the outside, from a single-threaded program holding one
copy of the data.
\end{jargon}

\begin{keyidea}{Why linearizability is the interesting property to check}
It is a property of the \emph{observable history} (what clients asked, and what
they were told), not of the implementation. That makes it checkable from the
outside, without trusting the code. If a system's history is linearizable, no
client could ever have detected that it was distributed. If it is not, some
client saw the seams: a value from the past, a write that vanished, a value
that never existed.
\end{keyidea}

\begin{figure}[H]
\centering
\resizebox{0.92\textwidth}{!}{
\begin{tikzpicture}[x=1cm, y=1cm]
\draw[->, thick, draw=inkblue] (0,0) -- (13.2,0) node[right, font=\small]{time};

% ok example
\node[font=\small\bfseries, text=okgreen, anchor=west] at (0,3.4) {Linearizable: a legal moment exists for each operation};
\draw[very thick, draw=accent] (0.6,2.7) -- (3.2,2.7);
\node[lbl, above] at (1.9,2.75) {C1: put city=lahore};
\filldraw[accent] (2.4,2.7) circle (2.2pt);
\node[lbl, below, text=okgreen] at (2.4,2.62) {takes effect};

\draw[very thick, draw=accent] (2.0,1.9) -- (4.6,1.9);
\node[lbl, above] at (3.3,1.95) {C2: get city $\rightarrow$ lahore};
\filldraw[accent] (3.4,1.9) circle (2.2pt);
\node[lbl, below, text=okgreen] at (3.4,1.82) {after the put: legal};

\draw[very thick, draw=accent] (5.4,2.7) -- (7.6,2.7);
\node[lbl, above] at (6.5,2.75) {C3: get city $\rightarrow$ lahore};
\filldraw[accent] (6.5,2.7) circle (2.2pt);

% bad example
\node[font=\small\bfseries, text=warnred, anchor=west] at (0,1.0) {Violation: no legal moment exists};
\draw[very thick, draw=warnred] (8.6,2.7) -- (10.2,2.7);
\node[lbl, above] at (9.4,2.75) {C1: put city=lahore};
\filldraw[warnred] (9.8,2.7) circle (2.2pt);
\draw[very thick, draw=warnred] (10.9,2.7) -- (12.8,2.7);
\node[lbl, above] at (11.85,2.75) {C2: get city $\rightarrow$ (nil)};
\node[lbl, below, text=warnred, align=center] at (11.85,2.62)
  {the put COMPLETED before\\this get STARTED, yet it\\was not seen: illegal};
\draw[dashed, draw=warnred] (10.2,2.5) -- (10.2,1.4);
\draw[dashed, draw=warnred] (10.9,2.5) -- (10.9,1.4);
\draw[<->, draw=warnred] (10.2,1.5) -- node[lbl, below, text=warnred]{no overlap: real-time order is forced} (10.9,1.5);
\end{tikzpicture}}
\caption{Left: operations overlap in real time, so the checker is free to pick
an order that explains every observed result. Right: two operations that do not
overlap have their order fixed by real time; if the results contradict that
order, no explanation exists and the history is a violation.}
\end{figure}

\subsection{The project's real thesis}

The README states the motivation directly: consensus code is easy to write and
nearly impossible to trust from example-based tests, because the bugs live in
message interleavings you will never hit on a laptop. So the project's central
bet, borrowed from the FoundationDB database's testing philosophy, is:

\begin{keyidea}{The thesis: make the whole distributed system a pure function of one integer}
Every source of randomness in a distributed system (what time it is, how long a
message takes, whether a message is lost, when a node crashes, when the network
splits, what the client does next) is replaced by draws from a single seeded
random number generator. A whole 12-second cluster lifetime under sustained
chaos then becomes a \emph{deterministic function of a seed}. That buys three
things at once: you can run thousands of universes in seconds because no real
time passes; a failure report is one integer rather than a shrug; and the
failing schedule replays exactly, forever, under a debugger.
\end{keyidea}

This only works if one condition holds, and the entire architecture is bent
around it: \textbf{the simulated system must be the real system}. If the
simulator tests a copy of the consensus logic, it tests nothing. The mechanism
that achieves this is the sans-IO core, section~\ref{sec:sansio}.

% =====================================================================
\section{The map: what is where}

\subsection{Directory structure}

\begin{figure}[H]
\centering
\begin{forest} dirtree
[raft-kv/
  [src/raftkv/
    [raft/ \quad{\rmfamily\itshape the consensus core: no sockets, no clocks, no asyncio}
      [node.py \quad{\rmfamily\itshape RaftNode, 576 lines: the algorithm itself}]
      [log.py \quad{\rmfamily\itshape snapshot-aware log index arithmetic}]
      [storage.py \quad{\rmfamily\itshape Storage protocol, FileStorage (fsync), MemoryStorage}]
      [messages.py \quad{\rmfamily\itshape the six RPC types, as dataclasses}]
      [env.py \quad{\rmfamily\itshape Env protocol: now / call\_later / send / rng}]
    ]
    [kv/machine.py \quad{\rmfamily\itshape the replicated state machine + session table}]
    [transport/tcp.py \quad{\rmfamily\itshape real asyncio TCP + AsyncioEnv (production)}]
    [sim/ \quad{\rmfamily\itshape the deterministic simulator (test-only)}
      [clock.py \quad{\rmfamily\itshape SimClock: one event heap, virtual time}]
      [network.py \quad{\rmfamily\itshape SimNetwork + SimEnv: loss, delay, partitions, skew}]
      [harness.py \quad{\rmfamily\itshape 439 lines: cluster, clients, nemesis, invariants}]
      [batch.py \quad{\rmfamily\itshape parallel seed sweep runner}]
    ]
    [checker/ \quad{\rmfamily\itshape the linearizability checker}
      [history.py \quad{\rmfamily\itshape Op and HistoryRecorder}]
      [linearizability.py \quad{\rmfamily\itshape Wing \& Gong search + Lowe memoization}]
    ]
    [server.py \quad{\rmfamily\itshape one TCP node: storage + machine + node + transport}]
    [client.py \quad{\rmfamily\itshape retrying client with leader hints and write dedup}]
    [cli/main.py \quad{\rmfamily\itshape the raftkv command}]
  ]
  [tests/ \quad{\rmfamily\itshape 57 tests across 10 files}]
  [docs/ \quad{\rmfamily\itshape ARCHITECTURE.md, TESTING.md, this document}]
  [demo.sh \quad{\rmfamily\itshape start cluster, kill -9 the leader, prove no data loss}]
  [pyproject.toml \quad{\rmfamily\itshape packaging; zero runtime dependencies}]
]
\end{forest}
\caption{The repository. The single most important boundary in the tree is that
\file{raft/} imports nothing from \file{sim/} or \file{transport/}: the arrow
of dependency points inward, always.}
\end{figure}

\subsection{The architecture in one picture}

\begin{figure}[H]
\centering
\resizebox{\textwidth}{!}{
\begin{tikzpicture}[node distance=10mm]

% core
\node[draw=inkblue, very thick, fill=softgrey, rounded corners=4pt,
      minimum width=54mm, minimum height=42mm] (core) {};
\node[font=\small\bfseries, text=inkblue, anchor=north] at ($(core.north)+(0,-2mm)$)
      {raft/ : the sans-IO core};
\node[box, minimum width=44mm] (rn) at ($(core.center)+(0,6mm)$) {RaftNode (node.py)};
\node[box, minimum width=20mm] (rl) at ($(core.center)+(-11mm,-6mm)$) {RaftLog};
\node[box, minimum width=20mm] (rm) at ($(core.center)+(12mm,-6mm)$) {Messages};
\node[lbl, anchor=south] at ($(core.south)+(0,2mm)$)
      {\itshape no import of asyncio, time, socket, or global random};

% two ports
\node[dec, above left=14mm and 6mm of core, minimum width=18mm] (env) {Env\\protocol};
\node[dec, above right=14mm and 6mm of core, minimum width=18mm] (sto) {Storage\\protocol};
\draw[flow] (rn) -- node[lbl, left]{now, call\_later,\\send, rng} (env);
\draw[flow] (rn) -- node[lbl, right]{load, append,\\truncate, snapshot} (sto);

% production side
\node[draw=okgreen, fill=white, rounded corners=4pt, minimum width=46mm,
      minimum height=30mm, above left=8mm and 4mm of env] (prod) {};
\node[font=\small\bfseries, text=okgreen, anchor=north] at ($(prod.north)+(0,-2mm)$) {PRODUCTION};
\node[extbox, minimum width=36mm] (aenv) at ($(prod.center)+(0,4mm)$) {AsyncioEnv + TcpTransport\\\scriptsize{}real clock, real sockets};
\node[extbox, minimum width=36mm] (fs) at ($(prod.center)+(0,-7mm)$) {FileStorage\\\scriptsize{}fsync, CRC, atomic rename};

% sim side
\node[draw=accent, fill=white, rounded corners=4pt, minimum width=46mm,
      minimum height=30mm, above right=8mm and 4mm of sto] (simb) {};
\node[font=\small\bfseries, text=accent, anchor=north] at ($(simb.north)+(0,-2mm)$) {SIMULATION};
\node[extbox, minimum width=36mm] (senv) at ($(simb.center)+(0,4mm)$) {SimEnv + SimNetwork + SimClock\\\scriptsize{}virtual time, seeded chaos};
\node[extbox, minimum width=36mm] (ms) at ($(simb.center)+(0,-7mm)$) {MemoryStorage\\\scriptsize{}a perfect disk};

\draw[dflow] (aenv) -- (env);
\draw[dflow] (fs) to[out=0,in=160] (sto);
\draw[dflow] (senv) to[out=180,in=20] (env);
\draw[dflow] (ms) -- (sto);

% consumers
\node[box, minimum width=26mm, below left=12mm and 2mm of core] (srv) {server.py\\\scriptsize{}RaftServer};
\node[box, minimum width=26mm, below right=12mm and 2mm of core] (har) {sim/harness.py\\\scriptsize{}SimCluster};
\draw[flow] (srv) -- (core);
\draw[flow] (har) -- (core);
\node[box, minimum width=24mm, below=8mm of srv] (cli) {cli/ + client.py};
\node[box, minimum width=24mm, below=8mm of har] (chk) {checker/};
\draw[flow] (cli) -- (srv);
\draw[flow] (har) -- node[lbl, right]{client history} (chk);
\node[store, minimum width=22mm, left=8mm of core] (kv) {kv/machine.py\\\scriptsize{}KVMachine};
\draw[flow] (rn) -- node[lbl, above, sloped]{apply(entry)} (kv);
\end{tikzpicture}}
\caption{The whole system. \code{RaftNode} sits in the middle and can only
touch the world through two narrow protocols, \code{Env} and \code{Storage}.
Production plugs real sockets and a real disk into those two sockets;
the simulator plugs in a virtual clock and a seeded chaos network. The
\emph{same} \code{RaftNode} object graph runs in both.}
\end{figure}

% =====================================================================
\section{The central design decision: sans-IO}
\label{sec:sansio}

\begin{jargon}{Sans-IO}
French for ``without IO''. A design style in which a protocol implementation
performs no input or output itself. It never opens a socket, never reads a
clock, never sleeps. Instead it is a pure-ish state machine: you feed it events
(``this message arrived'', ``this timer fired'', ``the local application wants
to propose a command''), and it returns or emits \emph{intentions} (``send this
message to n2'', ``call me back in 30 milliseconds''). Somebody else, the
so-called IO layer, actually performs those intentions.
\end{jargon}

\code{RaftNode} reacts to exactly three stimuli and produces effects through
exactly two protocols. That is the whole contract, and it fits in a box:

\begin{figure}[H]
\centering
\resizebox{0.98\textwidth}{!}{
\begin{tikzpicture}[node distance=8mm]
\node[box, minimum width=34mm, minimum height=30mm] (n) {\textbf{RaftNode}\\\scriptsize{}pure state machine};

\node[extbox, left=26mm of n, minimum width=34mm, yshift=13mm] (s1)
  {\code{handle\_message(src, msg)}\\\scriptsize{}a peer said something};
\node[extbox, left=26mm of n, minimum width=34mm] (s2)
  {timer callback\\\scriptsize{}scheduled earlier via call\_later};
\node[extbox, left=26mm of n, minimum width=34mm, yshift=-13mm] (s3)
  {\code{propose} / \code{read\_index}\\\code{add\_server} / \code{remove\_server}\\\scriptsize{}the local application asked};

\node[extbox, right=26mm of n, minimum width=34mm, yshift=13mm] (e1)
  {\code{env.send(dest, msg)}};
\node[extbox, right=26mm of n, minimum width=34mm] (e2)
  {\code{env.call\_later(d, fn)}};
\node[extbox, right=26mm of n, minimum width=34mm, yshift=-13mm] (e3)
  {\code{storage.save\_term\_vote}\\\code{storage.append\_entries} \dots};

\draw[flow] (s1) -- (n);
\draw[flow] (s2) -- (n);
\draw[flow] (s3) -- (n);
\draw[flow] (n) -- (e1);
\draw[flow] (n) -- (e2);
\draw[flow] (n) -- (e3);
\node[lbl, above=1mm of n] {IN: three stimuli \qquad OUT: two protocols};
\end{tikzpicture}}
\caption{The sans-IO contract of \code{RaftNode}. Nothing else crosses the
boundary.}
\end{figure}

The \code{Env} protocol is 6 meaningful lines, in \file{raft/env.py}:

\begin{lstlisting}[language=Python, caption={raft/env.py: the entire abstraction that makes the project testable}]
class Env(Protocol):
    rng: random.Random

    def now(self) -> float: ...
    def call_later(self, delay: float, fn: Callable[[], None]) -> TimerHandle: ...
    def send(self, dest: str, msg: Message) -> None: ...
\end{lstlisting}

\begin{keyidea}{Why this tiny file is the most important one in the repository}
Because \code{RaftNode} takes its randomness from \code{env.rng} rather than
the global \code{random} module, and its time from \code{env.now()} rather than
\code{time.time()}, the simulator can supply a seeded generator and a virtual
clock, and the node cannot tell the difference. That is what makes ``run 5000
chaos schedules in 90 seconds'' an ordinary function call. And crucially, it
means every bug the simulator finds is a bug in the code that runs in
production, because it \emph{is} the code that runs in production. There is no
test double of the algorithm anywhere in this repository.
\end{keyidea}

Note the direction of the dependency arrow. \code{Env} and \code{Storage} are
declared as \code{typing.Protocol}, which in Python means structural typing:
any object with the right method names satisfies them, with no inheritance and
no import. So \file{raft/} declares what it needs, and \file{sim/} and
\file{transport/} depend on \file{raft/}, never the reverse. This is dependency
inversion, and it is why \file{raft/node.py} contains no \code{if testing:}
branch anywhere.

% =====================================================================
\section{Raft: the algorithm, and this implementation of it}

Raft is defined by Diego Ongaro and John Ousterhout's 2014 paper ``In Search of
an Understandable Consensus Algorithm'', with log compaction and membership
changes taken from Ongaro's Stanford dissertation. This section explains the
algorithm from scratch and then shows exactly how \file{raft/node.py}
implements each piece.

\subsection{Terms: logical time}

\begin{jargon}{Term}
A number that only ever increases, acting as a logical clock. Each term begins
with an election. At most one node may be leader in a given term. Every message
carries its sender's term. The rule is unconditional and lives at the very top
of \code{handle\_message}: if you receive any message with a term greater than
your own, you immediately adopt that term and become a follower, whatever you
were doing. Terms let nodes detect stale information without any assumption
about real clocks being synchronized.
\end{jargon}

\begin{lstlisting}[language=Python, caption={raft/node.py: the term rule, applied before dispatch, to every message type}]
def handle_message(self, src: str, msg: Message) -> None:
    if msg.term > self.current_term:
        self._step_down(msg.term)
    if isinstance(msg, RequestVote):
        self._on_request_vote(msg)
    # ... dispatch to the other five handlers
\end{lstlisting}

\subsection{The three roles, and the state machine}

Every node is, at any instant, in exactly one of three roles.

\begin{itemize}
\item \textbf{Follower}: passive. Answers requests from leaders and candidates.
      Does nothing on its own except notice silence.
\item \textbf{Candidate}: campaigning. It got bored of the silence and is
      asking for votes.
\item \textbf{Leader}: in charge. It alone accepts client writes, it alone
      appends to the log, and it sends periodic heartbeats to suppress
      elections.
\end{itemize}

\begin{figure}[H]
\centering
\resizebox{0.99\textwidth}{!}{
\begin{tikzpicture}[font=\small]

\node[box, minimum width=32mm, minimum height=14mm] (f)
  {\textbf{FOLLOWER}\\\scriptsize{}passive, obeys the leader};
\node[box, minimum width=32mm, minimum height=14mm, right=44mm of f] (c)
  {\textbf{CANDIDATE}\\\scriptsize{}campaigning for votes};
\node[box, minimum width=32mm, minimum height=14mm, right=44mm of c] (l)
  {\textbf{LEADER}\\\scriptsize{}appends, replicates,\\\scriptsize{}heartbeats};

\node[lbl, left=12mm of f, align=center] (start) {start\\(also after a\\restart from disk)};
\draw[flow] (start) -- (f);

% self loops, on top
\draw[flow] (f) to[out=115, in=65, looseness=5.5]
  node[lbl, above, align=center, pos=0.5]{heartbeat received\\$\Rightarrow$ reset election timer} (f);
\draw[flow] (c) to[out=115, in=65, looseness=5.5]
  node[lbl, above, align=center, pos=0.5]{split vote:\\timeout again\\$\Rightarrow$ term += 1,\\campaign anew} (c);
\draw[flow] (l) to[out=115, in=65, looseness=5.5]
  node[lbl, above, align=center, pos=0.5]{heartbeat tick\\$\Rightarrow$ broadcast AppendEntries} (l);

% forward edges
\draw[flow] (f) -- node[lbl, below, align=center, pos=0.5]
  {\textbf{1. election timeout}\\silence for a random time in [T, 2T]\\$\Rightarrow$ term += 1, vote for self,\\fsync, send RequestVote} (c);
\draw[flow] (c) -- node[lbl, above, align=center, pos=0.5]
  {\textbf{2. votes from a majority}\\$\Rightarrow$ append a no-op entry\\of my own term,\\start heartbeat timer} (l);

% return edges, well separated vertically
\draw[flow] (c) to[out=235, in=305] node[lbl, below, align=center, pos=0.5]
  {\textbf{3. a current leader exists}\\(a valid AppendEntries arrives)\\or a higher term is seen} (f);

\draw[flow] (l) to[out=265, in=275, looseness=1.1] node[lbl, below, align=center, pos=0.5]
  {\textbf{4. a higher term is seen}, or I am removed from the config\\
   $\Rightarrow$ \code{\_step\_down}: cancel heartbeats, fail every pending read and write} (f);
\end{tikzpicture}}
\caption{The Raft role state machine as implemented in \file{raft/node.py}.
Note the asymmetry that makes Raft understandable: there is no
follower-to-leader edge and no leader-to-candidate edge. Power is only ever
taken through an election, and only ever surrendered by stepping down to
follower.}
\label{fig:statemachine}
\end{figure}

\begin{honest}{An implementation detail worth knowing: the timer never stops}
\code{\_on\_election\_timeout} re-arms itself \emph{before} checking whether it
should act, and returns early if the node is a leader or is not in its own
configuration. So the election timer keeps firing forever on every node, even
the leader, and even a removed node; it just does nothing. This is harmless and
simplifies the code (there is exactly one place that arms the election timer),
but it means a reader looking for ``where does the leader cancel its election
timer'' will not find it, because it does not.
\end{honest}

\subsection{Leader election, step by step}

\begin{jargon}{Election timeout}
How long a follower waits without hearing from a leader before it assumes the
leader is dead and starts campaigning. The critical detail is that it is
\textbf{randomized} per node and per attempt, drawn uniformly from
\code{[T, 2T]}. If it were a fixed constant, all followers would time out at the
same instant, all campaign at once, all split the vote, all fail, and repeat
forever. Randomization is what makes some node reliably go first and win.
\end{jargon}

\begin{lstlisting}[language=Python, caption={raft/node.py: randomized timeout, drawn from the injected rng, never the global one}]
def _reset_election_timer(self) -> None:
    if self._election_timer is not None:
        self._election_timer.cancel()
    delay = self.env.rng.uniform(self.election_timeout, 2 * self.election_timeout)
    self._election_timer = self.env.call_later(delay, self._on_election_timeout)
\end{lstlisting}

\begin{pseudocode}{Election, as implemented in \_start\_election and \_on\_request\_vote}
\Step{ON election timeout (follower or candidate, and I am in my own config):}
\Indent{role $\leftarrow$ CANDIDATE}
\Indent{currentTerm $\leftarrow$ currentTerm + 1}
\Indent{votedFor $\leftarrow$ me}
\Indent{\textbf{PERSIST} (term, votedFor) to storage \quad // before any message leaves}
\Indent{votes $\leftarrow$ \{me\}}
\Indent{IF votes $\geq$ majority THEN become leader \quad // single-node cluster}
\Indent{send RequestVote(term, me, lastLogIndex, lastLogTerm) to every peer}
\Step{}
\Step{ON RequestVote(m) received:}
\Indent{granted $\leftarrow$ false}
\Indent{IF m.term == currentTerm AND votedFor IN \{none, m.candidateId\}:}
\Indent{\quad upToDate $\leftarrow$ (m.lastLogTerm, m.lastLogIndex) $\geq$ (myLastLogTerm, myLastLogIndex)}
\Indent{\quad IF upToDate:}
\Indent{\quad\quad granted $\leftarrow$ true; votedFor $\leftarrow$ m.candidateId}
\Indent{\quad\quad \textbf{PERSIST} (term, votedFor); reset election timer}
\Indent{reply RequestVoteReply(currentTerm, me, granted)}
\Step{}
\Step{ON RequestVoteReply(m), if I am still a CANDIDATE in term m.term:}
\Indent{IF m.voteGranted AND m.voterId in my config: votes $\leftarrow$ votes $\cup$ \{m.voterId\}}
\Indent{IF |votes| $\geq$ majority: become leader}
\end{pseudocode}

Two subtleties in that pseudocode carry enormous weight.

\begin{keyidea}{Persist the vote before replying, or lose safety}
\code{\_persist\_term\_vote()} is called \emph{before} the reply is sent, both
when campaigning and when granting a vote. If a node granted a vote, crashed,
recovered with no memory of it, and granted a second vote for the same term to
a different candidate, two candidates could each collect a majority in one
term, and there would be two leaders. Durability of \code{votedFor} is not
bookkeeping; it is what makes ``at most one leader per term'' true across
crashes.
\end{keyidea}

\begin{keyidea}{The up-to-date check: elections are a filter, not a lottery}
A voter refuses any candidate whose log is less up to date than its own, where
``more up to date'' means: higher last term wins; if the last terms tie, the
longer log wins. In Python that comparison is literally a tuple comparison,
\code{(msg.last\_log\_term, msg.last\_log\_index) >= (self.log.last\_term,
self.log.last\_index)}. This is the paper's section 5.4.1, and it is what makes
the \emph{Leader Completeness} property true: since a winner needs a majority,
and any majority overlaps the majority that stored any committed entry, and
every voter refuses anyone missing entries it holds, a node missing a committed
entry can never win. Committed data therefore survives every leader change,
automatically, with no data transfer during the election itself.
\end{keyidea}

\begin{figure}[H]
\centering
\resizebox{0.94\textwidth}{!}{
\begin{tikzpicture}[x=1cm, y=1cm]
% lifelines
\foreach \x/\n in {0/n1, 4.6/n2, 9.2/n3} {
  \node[box, minimum width=18mm] at (\x,7.6) {\n};
  \draw[dashed, draw=linegrey, thick] (\x,7.2) -- (\x,-1.1);
}
% silence
\node[lbl, align=center, fill=softgrey, draw=linegrey] at (4.6,6.8)
  {the leader has died: no heartbeats reach anyone};

% n1 times out first
\node[lbl, anchor=east, align=right, text=warnred] at (-0.7,6.0) {election timeout fires\\(randomized: n1 first)};
\filldraw[warnred] (0,6.0) circle (2.5pt);
\node[lbl, anchor=east, align=right] at (-0.7,5.2) {term 1 $\rightarrow$ 2\\votedFor = n1\\\textbf{fsync before replying}};

\draw[flow] (0,4.5) -- node[lbl, above, pos=0.5]{RequestVote(term=2, lastIdx=5, lastTerm=1)} (4.6,4.5);
\draw[flow] (0,4.5) -- (9.2,4.5);

\node[lbl, anchor=north west, align=left] at (4.75,4.2) {term 1 $<$ 2: adopt term 2.\\my log is not newer:\\\textbf{grant}, fsync votedFor=n1};
\node[lbl, anchor=north west, align=left] at (9.35,4.2) {same: \textbf{grant}};

\draw[flow] (4.6,2.9) -- node[lbl, above, pos=0.5]{VoteGranted} (0,2.9);
\node[lbl, anchor=east, align=right, text=okgreen] at (-0.7,2.4) {votes = \{n1, n2\}\\2 $\geq$ majority(3) = 2\\\textbf{LEADER}};
\filldraw[okgreen] (0,2.9) circle (2.5pt);
\draw[flow] (9.2,2.2) -- node[lbl, below, pos=0.42]{VoteGranted (arrives late, already won: ignored)} (0,1.9);

\node[lbl, anchor=east, align=right] at (-0.7,1.3) {append no-op\\at index 6, term 2};

\draw[flow] (0,0.7) -- node[lbl, above, pos=0.5]{AppendEntries(term=2, entries=[6:noop])} (4.6,0.7);
\draw[flow] (0,0.7) -- (9.2,0.7);
\node[lbl, anchor=north west, align=left] at (4.75,0.4) {reset election timer;\\n1 is my leader};
\node[lbl, anchor=north west, align=left] at (9.35,0.4) {reset election timer};
\end{tikzpicture}}
\caption{One election, end to end. n1's randomized timer fires first, so it
campaigns alone and wins immediately; n3's vote arrives after the fact and is
ignored. The new leader's first act is to append a no-op entry, for the reason
explained in section~\ref{sec:commitrule}.}
\label{fig:election}
\end{figure}

\subsection{Log replication}

\begin{jargon}{Log entry}
One record in the ordered list. In \file{raft/messages.py} it is a frozen
dataclass with four fields: \code{index} (its 1-based position, never reused),
\code{term} (the term of the leader that created it), \code{kind}, and
\code{data}. There are exactly three kinds in this codebase:
\code{"command"} (a KV operation to run), \code{"config"} (a cluster membership
change), and \code{"noop"} (an empty entry a new leader appends to itself).
\end{jargon}

\begin{jargon}{Commit / committed}
An entry is \emph{committed} once the leader has stored it on a majority of the
cluster. Committed means permanent: it will be present in every future leader's
log and will eventually be applied on every node. Uncommitted entries can still
be thrown away. \code{commit\_index} is the highest index known committed;
\code{last\_applied} is the highest index actually executed against the state
machine. The gap between them is entries that are decided but not yet run.
\end{jargon}

The leader tracks two numbers per follower, and understanding the difference is
most of understanding replication:

\begin{center}
\begin{tabular}{lll}
\toprule
\textbf{Field} & \textbf{Meaning} & \textbf{Nature} \\
\midrule
\code{next\_index[p]} & the next index I will \emph{try} to send p & optimistic guess \\
\code{match\_index[p]} & the highest index I \emph{know} p has stored & proven fact \\
\bottomrule
\end{tabular}
\end{center}

On becoming leader, \code{next\_index} is set optimistically to
\code{last\_index + 1} for everyone and \code{match\_index} to 0 for everyone:
``I assume you are caught up, and I have proof of nothing.'' Reality then
corrects the guess.

\begin{pseudocode}{AppendEntries on the follower (\_on\_append\_entries), the consistency check}
\Step{ON AppendEntries(m):}
\Indent{IF m.term $<$ currentTerm: reply(success=false) \quad // stale leader, ignore}
\Indent{IF role != FOLLOWER: step down \quad // a valid leader exists for this term}
\Indent{leaderHint $\leftarrow$ m.leaderId; reset election timer}
\Step{}
\Indent{// THE LOG MATCHING CHECK: do we agree on the entry just before these?}
\Indent{IF term\_at(m.prevLogIndex) != m.prevLogTerm:}
\Indent{\quad // We disagree. Tell the leader WHERE, so it can back up fast.}
\Indent{\quad IF we have no entry there at all: conflict $\leftarrow$ myLastIndex + 1}
\Indent{\quad ELSE: conflict $\leftarrow$ first index of the conflicting term}
\Indent{\quad reply(success=false, conflictIndex=conflict); RETURN}
\Step{}
\Indent{// We agree up to prevLogIndex. Merge the new entries.}
\Indent{FOR each entry e in m.entries:}
\Indent{\quad IF we have nothing at e.index: append it}
\Indent{\quad ELSE IF our term at e.index != e.term:}
\Indent{\quad\quad TRUNCATE our log from e.index onward \quad // ours was wrong; leader wins}
\Indent{\quad\quad append e}
\Indent{\quad // ELSE: identical entry already present, skip}
\Indent{persist appended entries (fsync); refresh config if a config entry moved}
\Indent{IF m.leaderCommit $>$ commitIndex: commitIndex $\leftarrow$ min(m.leaderCommit, match)}
\Indent{reply(success=true, matchIndex=m.prevLogIndex + len(m.entries), seq=m.seq)}
\end{pseudocode}

\begin{keyidea}{The Log Matching Property, and why one check suffices}
The check ``do we agree on the single entry immediately before the ones I am
sending'' looks far too weak to guarantee that two whole logs agree. It is not,
because of an induction: entries are only ever appended by a leader, indexes
are never reused, and only one leader exists per term. So if two logs contain an
entry with the same index \emph{and} the same term, that entry was created by
the same leader in the same term, and therefore all entries before it are also
identical. Agreement on one entry proves agreement on the entire prefix. That is
why AppendEntries carries \code{prev\_log\_index} and \code{prev\_log\_term}
at all, and why the follower's reply is trustworthy.
\end{keyidea}

\begin{figure}[H]
\centering
\resizebox{\textwidth}{!}{
\begin{tikzpicture}[x=1cm, y=1cm]
\foreach \x/\n/\r in {0/n1/{leader}, 6.2/n2/{follower}, 12.4/n3/{follower}} {
  \node[box, minimum width=20mm] at (\x,9.4) {\n\\\scriptsize\r};
  \draw[dashed, draw=linegrey, thick] (\x,8.9) -- (\x,-1.9);
}
\node[box, minimum width=18mm] at (-3.8,9.4) {Client};
\draw[dashed, draw=linegrey, thick] (-3.8,8.9) -- (-3.8,-1.9);

\draw[flow] (-3.8,8.3) -- node[lbl, above]{put city=lahore} (0,8.3);
\node[lbl, anchor=north west, align=left, text width=5.4cm] at (0.3,8.0)
  {\code{propose()}: append Entry(idx=6, term=2, command)\\
   to my own log, \textbf{fsync it}, then broadcast};

\draw[flow] (0,6.2) -- node[lbl, above, pos=0.5]
  {AppendEntries(prevIdx=5, prevTerm=2, entries=[6], leaderCommit=5)} (6.2,6.2);
\node[lbl, anchor=north west, align=left, text width=4.6cm] at (6.5,6.1)
  {my term at index 5 is 2:\\\textbf{matches}. Append entry 6,\\\textbf{fsync}};

\draw[flow] (0,4.7) -- node[lbl, above, pos=0.5]{AppendEntries(\dots), the same entry to n3} (9.15,4.7);
\node[font=\Large, text=warnred] at (9.4,4.7) {$\times$};
\node[lbl, anchor=north, align=center, text width=3.4cm, text=warnred] at (9.4,4.4)
  {this message is lost\\by the network};

\draw[flow] (6.2,3.4) -- node[lbl, above, pos=0.5]{Reply(success=true, matchIndex=6)} (0,3.4);

\node[lbl, anchor=north west, align=left, text width=7.4cm, text=okgreen] at (0.3,3.1)
  {matchIndex = \{n2:6, n3:0\}, plus my own 6.\\
   sorted = [0, 6, 6]; the majority-th value is 6.\\
   Entry 6 is of MY term (2), so it may commit.\\
   \textbf{commitIndex = 6}. Apply to the state machine.};

\draw[flow] (0,0.9) -- node[lbl, above, pos=0.5]{reply: OK} (-3.8,0.9);
\node[lbl, anchor=north west, align=left, text width=7.4cm] at (0.3,0.6)
  {n3 is still behind, but the client already has its answer:\\a majority is enough.};

\draw[flow] (0,-1.0) -- node[lbl, above, pos=0.5]
  {next heartbeat: AppendEntries(entries=[6], leaderCommit=6)} (12.4,-1.0);
\node[lbl, anchor=north west, align=left, text width=3.2cm] at (12.7,-1.2) {n3 catches up here};
\end{tikzpicture}}
\caption{The replication path for one write. The client is answered as soon as a
\emph{majority} has the entry stored durably; the straggler catches up on the
next heartbeat. Losing n3's message costs nothing, because AppendEntries is
idempotent and heartbeats retry forever.}
\label{fig:replication}
\end{figure}

\subsubsection{Fast backoff on conflict}

When a follower rejects an AppendEntries, the naive repair is for the leader to
decrement \code{next\_index} by one and try again, costing one full round trip
per missing entry. A follower 500 entries behind then needs 500 round trips.
Raft's optimization, implemented here, is for the follower to return a
\emph{conflict hint}: the first index of the conflicting term. The leader jumps
straight there.

\begin{lstlisting}[language=Python, caption={raft/node.py: the follower computes the hint by scanning back over its own conflicting term}]
our_prev_term = self.log.term_at(prev_index)
if our_prev_term is None or our_prev_term != prev_term:
    if our_prev_term is None:
        conflict = self.log.last_index + 1        # our log is simply too short
    else:
        conflict = prev_index                     # walk back over the whole bad term
        while conflict > self.log.snapshot_index + 1 and \
                self.log.term_at(conflict - 1) == our_prev_term:
            conflict -= 1
    self.env.send(msg.leader_id, AppendEntriesReply(
        term=self.current_term, follower_id=self.id, success=False,
        match_index=0, conflict_index=conflict, seq=msg.seq))
    return
\end{lstlisting}

\begin{lstlisting}[language=Python, caption={raft/node.py: and the leader consumes it, but never trusts it upward}]
else:
    # Fast backoff using the follower's conflict hint.
    self.next_index[p] = max(1, min(self.next_index[p] - 1, msg.conflict_index)
                             if msg.conflict_index else self.next_index[p] - 1)
    self._send_append(p)
\end{lstlisting}

The \code{min(next\_index - 1, conflict\_index)} is a defensive detail worth
noticing in an interview: the hint can only ever move \code{next\_index}
\emph{down}, never up. A hint from a confused or reordered reply can therefore
cost an extra round trip but can never skip entries, which would be a
correctness bug rather than a performance one. The \code{max(1, ...)} floors it
at the first real index.

\subsection{The commit rule, and the no-op}
\label{sec:commitrule}

This is the single subtlest rule in Raft, and the README singles it out as one
of the two places ``where a plausible shortcut is a real linearizability bug''.

\begin{lstlisting}[language=Python, caption={raft/node.py: \_maybe\_advance\_commit, all seven lines of it}]
def _maybe_advance_commit(self) -> None:
    if self.role is not Role.LEADER:
        return
    matches = sorted(
        (self.match_index[n] if n != self.id else self.log.last_index)
        for n in self.config
    )
    if not matches:
        return
    candidate = matches[len(matches) - self._majority()]
    if candidate > self.commit_index and self.log.term_at(candidate) == self.current_term:
        self._set_commit_index(candidate)
\end{lstlisting}

The first half is a neat trick: sort every node's match index (counting the
leader's own last index for itself) and take the element at position
\code{len - majority}. For 3 nodes with matches \code{[0, 6, 6]}, that is index
\code{3 - 2 = 1}, giving 6: the highest index that a majority has reached. It is
the median, generalized.

The second half is the subtle part: \code{self.log.term\_at(candidate) ==
self.current\_term}. A leader may \textbf{not} commit an entry from an earlier
term merely because it is now replicated on a majority.

\begin{keyidea}{Why counting replicas of an old entry is unsafe (the paper's Figure 8)}
Suppose an entry from term 2 sits on a majority of nodes but was never
committed by term 2's leader (which crashed before it could count). A new
leader in term 4 sees it on a majority and commits it. But a \emph{different}
node, whose log happens to be more up to date on other grounds, can still win a
later election, and its log does not contain that entry. It would then
overwrite it. An entry that was reported committed would vanish. Replication
count alone is therefore not proof of permanence for old entries.

The fix, and this implementation's version of it: a leader may only directly
commit entries \textbf{from its own term}. To make progress on the backlog, a
fresh leader immediately appends a \code{noop} entry of its own term
(\code{\_become\_leader}, line 212). When that no-op commits by the ordinary
majority rule, every entry before it commits \emph{transitively}, by the Log
Matching Property, safely. The no-op is not a nicety; without it a leader that
receives no writes could never commit its predecessor's backlog at all.
\end{keyidea}

\begin{figure}[H]
\centering
\resizebox{0.98\textwidth}{!}{
\begin{tikzpicture}[x=0.92cm, y=0.92cm]
\tikzset{cell/.style={draw=inkblue, minimum width=8mm, minimum height=8mm, font=\small, fill=white}}
\tikzset{cellold/.style={cell, fill=softgrey}}
\tikzset{cellnew/.style={cell, fill=okgreen!18, draw=okgreen}}

\node[lbl, anchor=east, align=right] at (-0.4,3.0) {\textbf{n1} (new leader, term 2)};
\foreach \i/\t in {0/1, 1/1, 2/2} {
  \node[cellold] at (\i,3.0) {\t};
}
\node[cellnew] at (3,3.0) {2};
\node[lbl, below] at (3,2.5) {noop};

\node[lbl, anchor=east, align=right] at (-0.4,1.8) {\textbf{n2}};
\foreach \i/\t in {0/1, 1/1} { \node[cellold] at (\i,1.8) {\t}; }
\node[cell, draw=linegrey, fill=white, text=linegrey] at (2,1.8) {?};
\node[cell, draw=linegrey, fill=white, text=linegrey] at (3,1.8) {?};

\node[lbl, anchor=east, align=right] at (-0.4,0.6) {\textbf{n3}};
\foreach \i/\t in {0/1, 1/1} { \node[cellold] at (\i,0.6) {\t}; }
\node[cell, draw=linegrey, fill=white, text=linegrey] at (2,0.6) {?};
\node[cell, draw=linegrey, fill=white, text=linegrey] at (3,0.6) {?};

\foreach \i in {0,1,2,3} { \node[lbl, above] at (\i,3.5) {idx \the\numexpr\i+1\relax}; }

\node[lbl, anchor=west, align=left, text width=7.6cm, text=warnred] at (4.4,3.1)
  {\textbf{FORBIDDEN}: n1 may not commit index 3 (term 1... here term 2 but
   created by the PREVIOUS leader) just because a majority stores it. Old-term
   entries are not committable by replica count.};
\node[lbl, anchor=west, align=left, text width=7.6cm, text=okgreen] at (4.4,1.2)
  {\textbf{ALLOWED}: once the term-2 noop at index 4 reaches a majority, index 4
   commits by the ordinary rule, and indexes 1 to 3 commit transitively with it,
   because any log holding entry 4 must hold its whole prefix.};
\draw[flow, draw=okgreen] (3,2.4) to[out=200, in=90] (2.0,1.0);
\end{tikzpicture}}
\caption{The section 5.4.2 commit restriction. The no-op entry is the vehicle
that safely drags the previous leader's backlog into the committed region.}
\end{figure}

The test \code{test\_commit\_only\_with\_current\_term\_entry} pins exactly
this: it force-installs a leader with two term-1 entries, asserts
\code{commit\_index == 0} immediately (nothing commits before replication),
runs, and asserts \code{commit\_index == 3}: the no-op at index 3 committed and
pulled indexes 1 and 2 with it.

\subsection{Persistence and crash safety}

\begin{jargon}{fsync}
An operating system call that does not return until the data you wrote has
actually reached durable storage. An ordinary \code{write()} only copies data
into a kernel buffer in memory; if the power fails a moment later, that data
never existed. Without \code{fsync}, all of Raft's durability reasoning is
fiction.
\end{jargon}

\code{FileStorage} keeps three files per node and uses two different, carefully
chosen durability techniques:

\begin{figure}[H]
\centering
\resizebox{\textwidth}{!}{
\begin{tikzpicture}[node distance=8mm]
\node[store, minimum width=32mm, minimum height=16mm] (h) {\code{hardstate.json}\\\scriptsize{}currentTerm, votedFor};
\node[store, minimum width=32mm, minimum height=16mm, right=16mm of h] (l) {\code{log.bin}\\\scriptsize{}append-only records};
\node[store, minimum width=32mm, minimum height=16mm, right=16mm of l] (s) {\code{snapshot.json}\\\scriptsize{}state machine image};

\node[extbox, below=9mm of h, minimum width=42mm, align=left] (h2)
  {\scriptsize \textbf{Technique: atomic replace}\\
   \scriptsize write .tmp $\rightarrow$ flush $\rightarrow$ fsync file\\
   \scriptsize $\rightarrow$ rename() $\rightarrow$ fsync the DIRECTORY\\
   \scriptsize rename is atomic on POSIX: a reader\\
   \scriptsize sees the whole old file or the whole new one};
\node[extbox, below=9mm of l, minimum width=42mm, align=left] (l2)
  {\scriptsize \textbf{Technique: scan and truncate}\\
   \scriptsize each record = [4B length][4B CRC32][JSON]\\
   \scriptsize append $\rightarrow$ flush $\rightarrow$ fsync, per batch\\
   \scriptsize on load: scan; first short or bad-CRC\\
   \scriptsize record = torn tail $\rightarrow$ truncate from there};
\node[extbox, below=9mm of s, minimum width=42mm, align=left] (s2)
  {\scriptsize \textbf{Technique: atomic replace, then rewrite}\\
   \scriptsize snapshot.json lands atomically, THEN\\
   \scriptsize log.bin is rewritten with only the suffix.\\
   \scriptsize Crash in between? On load, entries at or\\
   \scriptsize below snapshot.last\_index are dropped.};
\draw[flow] (h) -- (h2);
\draw[flow] (l) -- (l2);
\draw[flow] (s) -- (s2);
\end{tikzpicture}}
\caption{The three on-disk structures and their recovery strategies. Each is
either atomic-rename recoverable or scan-and-truncate recoverable, which is the
project's stated one-line discipline for crash safety.}
\end{figure}

\begin{keyidea}{Why a torn tail is always safe to discard}
A crash mid-append can leave half a record on disk. Recovery finds the bad
CRC32 checksum and truncates from that point. This throws away data, which
sounds alarming, and is provably fine: Raft never acknowledges anything before
\code{fsync} returns, so a record that was still being written was, by
definition, never acknowledged to anybody. Nobody was told it existed, so
nobody can be surprised that it does not. Discipline (never acknowledge before
fsync) is what makes a destructive recovery correct.
\end{keyidea}

The scan is genuinely as small as the README claims. This is all of it:

\begin{lstlisting}[language=Python, caption={raft/storage.py: torn-tail recovery, the whole mechanism}]
while pos + _REC_HEADER.size <= len(raw):
    length, crc = _REC_HEADER.unpack_from(raw, pos)
    start = pos + _REC_HEADER.size
    payload = raw[start:start + length]
    if len(payload) < length or zlib.crc32(payload) != crc:
        break  # torn tail from a crash mid-write
    entries.append(Entry.from_dict(json.loads(payload)))
    pos = start + length
    good_end = pos
if good_end < len(raw):
    # Discard the torn tail so future appends start from a clean point.
    self._log_fh.close()
    with open(self._logf, "r+b") as f:
        f.truncate(good_end)
        f.flush()
        os.fsync(f.fileno())
    self._log_fh = open(self._logf, "ab")
\end{lstlisting}

Note the detail in \code{\_atomic\_write}: after \code{os.rename}, the code
opens the containing \emph{directory} and fsyncs that too. Renaming is atomic,
but the directory entry recording the rename is itself buffered; without the
directory fsync, a power cut can leave the file at its old name. This is a
classic mistake, and it is not made here.

Two tests drive these paths directly rather than by inspection:

\begin{itemize}
\item \code{test\_crash\_mid\_write\_discards\_torn\_tail}\\
      writes a valid record header promising 500 bytes and then supplies 11,
      then asserts both that the surviving entries are exactly
      \code{[1, 2, 3]} and that the file was physically truncated back to its
      previous size.
\item \code{test\_corrupt\_record\_stops\_scan}\\
      flips a single bit inside the last record's payload and asserts the CRC
      rejects it.
\end{itemize}

\subsection{Snapshots and log compaction}

An append-only log grows forever. Compaction replaces a prefix of the log with
a single image of the state that prefix produced.

\begin{jargon}{Snapshot}
A serialized copy of the entire state machine at a particular log index, plus
the metadata needed to keep the log arithmetic consistent:
\code{last\_included\_index}, \code{last\_included\_term}, and the cluster
config as of that point. Here the payload is \code{KVMachine.snapshot()}: a
JSON blob of the data dictionary \emph{and} the session table.
\end{jargon}

\begin{figure}[H]
\centering
\resizebox{0.95\textwidth}{!}{
\begin{tikzpicture}[x=0.85cm, y=0.85cm]
\tikzset{cell/.style={draw=inkblue, minimum width=7.5mm, minimum height=7.5mm, font=\scriptsize, fill=white}}
\node[lbl, anchor=east] at (-0.5,2.2) {before:};
\foreach \i in {0,...,9} { \node[cell, fill=softgrey] at (\i,2.2) {\the\numexpr\i+1\relax}; }
\foreach \i in {10,11} { \node[cell] at (\i,2.2) {\the\numexpr\i+1\relax}; }
\draw[decorate, decoration={brace, amplitude=4pt}, draw=accent]
  (-0.42,2.65) -- (9.42,2.65) node[midway, above=3pt, lbl, text=accent]{applied prefix: replaceable};

\node[lbl, anchor=east] at (-0.5,0.2) {after:};
\node[store, minimum width=42mm, minimum height=8mm] at (2.3,0.2)
  {\scriptsize snapshot: lastIndex=10, lastTerm=6, config, KV data + sessions};
\foreach \i in {10,11} { \node[cell] at (\i,0.2) {\the\numexpr\i+1\relax}; }
\draw[flow] (4.7,1.7) -- (4.7,0.75);
\node[lbl, anchor=west, align=left, text width=5.4cm] at (12.4,1.2)
  {\code{log.compact\_to(10, 6)}\\then\\\code{storage.save\_snapshot(...)}};
\end{tikzpicture}}
\caption{Compaction. \code{RaftLog} keeps \code{snapshot\_index} as an offset,
so index arithmetic on the surviving entries stays correct: \code{entry(i)}
resolves to \code{entries[i - snapshot\_index - 1]}.}
\end{figure}

The trigger is in \code{\_maybe\_snapshot}, called after every apply:

\begin{lstlisting}[language=Python]
if self.last_applied - self.log.snapshot_index < self.snapshot_threshold:
    return
\end{lstlisting}

Above the threshold (1000 entries in production, 40 in the simulator, 10 in the
snapshot tests) it takes an image and compacts.

The complication compaction creates is that a leader may no longer possess the
entries a lagging follower needs, because it deleted them. \code{\_send\_append}
detects this in its first two lines and switches protocol:

\begin{lstlisting}[language=Python, caption={raft/node.py: the fork between replication and snapshot transfer}]
def _send_append(self, p: str) -> None:
    ni = self.next_index[p]
    if ni <= self.log.snapshot_index:
        snap = self._load_snapshot()
        self.env.send(p, InstallSnapshot(
            term=self.current_term, leader_id=self.id,
            last_included_index=snap.last_index,
            last_included_term=snap.last_term,
            config=tuple(snap.config), data=snap.data))
        return
    prev_term = self.log.term_at(ni - 1)
    assert prev_term is not None
    entries = tuple(self.log.slice_from(ni))
    self.env.send(p, AppendEntries(...))
\end{lstlisting}

On the receiving side, \code{\_on\_install\_snapshot} contains a nice piece of
care. If the follower's own log already contains an entry at
\code{last\_included\_index} \emph{with the matching term}, the snapshot is
merely a prefix of what it already has, so it compacts and keeps its suffix.
Only if the terms disagree does it discard its log wholesale:

\begin{lstlisting}[language=Python, caption={raft/node.py: keep the suffix when the snapshot is only a prefix of what we hold}]
if self.log.term_at(msg.last_included_index) == msg.last_included_term:
    # Our log already extends past the snapshot; keep the suffix.
    self.log.compact_to(msg.last_included_index, msg.last_included_term)
else:
    self.log = RaftLog(msg.last_included_index, msg.last_included_term)
\end{lstlisting}

\begin{honest}{Snapshots are read from disk and shipped whole, on every send}
\code{\_load\_snapshot()} is literally \code{return self.storage.load().snapshot}.
For \code{FileStorage} that re-reads and re-parses \code{snapshot.json} from
disk \textbf{every single time} a snapshot is sent, and \code{load()} also
re-reads and re-parses the entire \code{log.bin} while it is there, discarding
the result. It is then sent as one message with no chunking, so a large state
machine means a single enormous frame (bounded only by \code{MAX\_FRAME}, 64
MB). At three nodes and a demo-sized dataset this is invisible; at any real
scale it is the first thing that would fall over. The README lists the chunked
transfer and cached handle as known future work, which is honest, but the
severity is easy to understate: the cost here is not just ``re-reads the
snapshot'', it is ``re-reads and re-parses the entire log too''.
\end{honest}

\subsection{Membership changes}

Changing the cluster from \{n1, n2, n3\} to \{n1, n2, n3, n4\} is dangerous for
a reason that is easy to miss: during the switch, different nodes believe in
different configurations, and if the old and new majorities do not overlap, two
leaders can be elected simultaneously without either breaking any rule.

\begin{keyidea}{Why single-server changes are safe without a two-phase protocol}
Add or remove exactly \textbf{one} server per config entry. Any two
configurations differing by one server must share a common majority: a majority
of \{n1,n2,n3\} is 2, a majority of \{n1,n2,n3,n4\} is 3, and any 2-of-3 and
any 3-of-4 drawn from these overlapping sets must intersect. So the disjoint
quorums that cause dual leadership cannot form, and no special ``joint'' state
is needed. Ongaro's dissertation's alternative, \emph{joint consensus}, allows
arbitrary jumps but costs a two-phase protocol and a state where both old and
new majorities must agree. \file{docs/ARCHITECTURE.md} states the tradeoff and
the choice explicitly.
\end{keyidea}

The implementation is a plain log entry of kind \code{config} plus exactly two
extra rules:

\begin{enumerate}
\item \textbf{One at a time.} \code{\_config\_change\_pending()} scans the
      uncommitted tail for any config entry and raises
      \code{ConfigChangeInProgress} if one is there.
\item \textbf{Effective on append, not on commit.} \code{\_refresh\_config()}
      runs the moment a config entry is appended or truncated. This is
      counter-intuitive (acting on data that might still be rolled back) but is
      the dissertation's rule, and it is necessary: a node must count itself
      out of a cluster it has been removed from immediately, or it can disrupt
      the very commit that removes it.
\end{enumerate}

Leader self-removal is handled at the end of \code{\_apply\_committed}: the
leader keeps leading until the config entry that removes it \emph{commits}, and
only then steps down, so it can shepherd its own removal to completion.

\begin{lstlisting}[language=Python, caption={raft/node.py: a leader that has been removed from the config steps down, but only on apply}]
if entry.kind == "config" and self.role is Role.LEADER and \
        self.id not in tuple(entry.data["nodes"]):
    self._step_down(self.current_term)
\end{lstlisting}

A joining node is started with an empty \code{initial\_config}. Because
\code{\_on\_election\_timeout} returns early when \code{self.id not in
self.config}, and an empty config contains nothing, a joiner never campaigns.
It sits silent until the leader's AppendEntries or InstallSnapshot teaches it
the cluster it belongs to. This is elegant: one guard clause implements
``joining nodes must not disrupt the cluster''.

\begin{honest}{A real dead end: \_config\_from's entries\_last\_index filter}
\code{\_config\_from(entries\_last\_index)} filters entries above a bound, but
the only caller that passes a bound is \code{\_maybe\_snapshot}, and it passes
\code{self.last\_applied}, which by construction is a prefix of the log at that
moment. The other caller, \code{\_latest\_config}, passes \code{None}. The
filter is therefore correct but currently near-vacuous. Not a bug; worth knowing
before an interviewer asks what it defends against.
\end{honest}

\begin{honest}{An inconsistency between two nearly identical lines}
\code{\_become\_leader} initializes \code{next\_index} to
\code{self.log.last\_index + 1}. \code{\_change\_config} initializes newly
appearing peers to \code{self.log.last\_index} (no \code{+ 1}), while
\code{\_refresh\_config} uses \code{last\_index + 1} again. The \code{+ 1} is
the paper's value. The odd one out in \code{\_change\_config} is off by one and
means the leader will send the new server one entry it need not have sent (the
config entry itself, which it does not have anyway, so the follower rejects and
the conflict hint repairs it in one extra round trip). It is self-correcting
and therefore invisible in the tests, but the three lines should agree, and an
interviewer reading these three functions side by side will notice.
\end{honest}

% =====================================================================
\section{Linearizable reads: the ReadIndex protocol}

Writes go through the log, so they are safe by construction. Reads are the trap.

\begin{keyidea}{Why reading from the leader's memory is wrong}
``I am the leader, so I will just answer from my local copy'' is unsafe,
because \textbf{a leader cannot know it is still the leader}. Suppose n1 is
partitioned away from n2 and n3. Those two elect n3 in a new term and accept
writes. n1 has heard nothing, so it still believes it is leader, its local data
looks perfectly fine to it, and it will happily serve a value that is now
minutes out of date. Nothing about its own state betrays it. Leadership is not
a local fact.
\end{keyidea}

The fix is to make the leader \emph{prove} its currency to a majority before
answering, which is the ReadIndex protocol:

\begin{pseudocode}{ReadIndex (\code{read\_index} + \code{\_check\_pending\_reads})}
\Step{ON read\_index(callback):}
\Indent{IF not leader: callback(None); RETURN \quad // caller must redirect}
\Indent{// 1. A fresh leader must first commit an entry of its OWN term.}
\Indent{readIdx $\leftarrow$ commitIndex IF an entry of my term has committed ELSE none}
\Indent{// 2. Record the current heartbeat round number and broadcast it.}
\Indent{hbSeq $\leftarrow$ hbSeq + 1}
\Indent{pendingReads.append(PendingRead(seq=hbSeq, readIndex=readIdx, cb=callback))}
\Indent{send AppendEntries(seq=hbSeq) to every peer}
\Step{}
\Step{ON any AppendEntriesReply(seq=s) from peer p:  ackSeq[p] $\leftarrow$ max(ackSeq[p], s)}
\Step{}
\Step{CHECK pending reads (after every reply, every commit, every broadcast):}
\Indent{FOR each pending read r with r.readIndex != none:}
\Indent{\quad acks $\leftarrow$ count of nodes n where n == me OR ackSeq[n] $\geq$ r.seq}
\Indent{\quad IF acks $\geq$ majority AND lastApplied $\geq$ r.readIndex:}
\Indent{\quad\quad remove r; r.cb(r.readIndex) \quad // NOW reading local state is safe}
\Step{}
\Step{ON step down (any cause): FOR each pending read r: r.cb(none) \quad // fail them all}
\end{pseudocode}

\begin{figure}[H]
\centering
\resizebox{\textwidth}{!}{
\begin{tikzpicture}[x=1cm, y=1cm]
% partition backdrop
\fill[warnred!7, rounded corners=3pt] (-3.1,-1.6) rectangle (2.0,8.5);
\fill[okgreen!7, rounded corners=3pt] (2.6,-1.6) rectangle (12.8,8.5);
\node[lbl, text=warnred, anchor=north west, align=left] at (-3.0,8.35)
  {\textbf{MINORITY SIDE} (1 of 3)\\cannot act};
\node[lbl, text=okgreen, anchor=north east, align=right] at (12.7,8.35)
  {\textbf{MAJORITY SIDE} (2 of 3): makes progress};

\draw[ultra thick, dashed, draw=warnred] (2.3,-1.6) -- (2.3,7.6);
\node[lbl, text=warnred, rotate=90, fill=white, inner sep=3pt] at (2.3,8.1)
  {\textbf{NETWORK PARTITION}};

% lifelines
\foreach \x/\n/\r in {0/n1/{was leader, term 5}, 4.6/n2/{follower}, 9.2/n3/{follower}} {
  \node[box, minimum width=22mm] at (\x,7.0) {\n\\\scriptsize\r};
  \draw[dashed, draw=linegrey, thick] (\x,6.5) -- (\x,-1.3);
}
\node[box, minimum width=20mm] at (-2.4,7.0) {Client};
\draw[dashed, draw=linegrey, thick] (-2.4,6.5) -- (-2.4,-1.3);

\draw[flow] (-2.4,5.9) -- node[lbl, above]{get city} (0,5.7);
\node[lbl, anchor=west, align=left, text width=2.4cm] at (0.3,5.2)
  {I am leader.\\readIdx = 42\\hbSeq = 7};

\draw[flow, draw=warnred] (0,4.2) -- node[lbl, above, pos=0.42]{AppendEntries(seq=7)} (2.05,4.2);
\node[font=\Large, text=warnred] at (2.3,4.2) {$\times$};
\node[lbl, anchor=west, text=warnred] at (2.6,4.2) {dropped by the partition};

\node[lbl, anchor=west, align=left, text width=7.0cm, text=okgreen] at (4.9,5.9)
  {\textbf{Meanwhile, on the majority side:} n2 and n3 time out,\\
   elect n3 in term 6, and commit new writes.\\
   The value of city is now 'karachi'.};
\draw[flow, draw=okgreen] (9.2,4.8) -- node[lbl, above]{AppendEntries(term=6)} (4.6,4.8);

\node[lbl, anchor=west, align=left, text width=2.6cm, text=warnred] at (0.3,2.7)
  {acks for seq 7\\= 1 (just me)\\1 $<$ majority(2)};
\filldraw[warnred] (0,2.7) circle (2.5pt);
\node[lbl, anchor=west, align=left, text width=2.6cm, text=warnred] at (0.3,1.4)
  {\textbf{the read never}\\\textbf{completes}};

\node[lbl, anchor=north west, align=left, text width=9.4cm] at (3.0,1.2)
  {When the partition heals, n1 sees term 6 $>$ 5 and calls \code{\_step\_down},
   which fails every pending read with \code{cb(None)}. The client is told
   \code{not\_leader}, retries at n3, and reads 'karachi'.\\
   \textbf{The stale value was never served.}};
\end{tikzpicture}}
\caption{A network partition, and why ReadIndex is safe. The deposed leader is
not merely unlikely to answer; it is structurally unable to, because the acks it
needs cannot cross the partition. Its own state gives it no clue that anything
is wrong, which is exactly why local state must never be trusted.}
\label{fig:partition}
\end{figure}

\begin{keyidea}{The second precondition: a fresh leader must commit its own term first}
The README calls this out as the bug the chaos harness found. A newly elected
leader's \code{commit\_index} is a stale local variable: it reflects what
\emph{it} knew before the election, and \emph{may be lower than what is actually
committed cluster-wide}, because commit index is not transferred during
elections. Capturing \code{readIndex = commitIndex} at that moment could
therefore serve a read that misses a committed write, which is a real
linearizability violation.

The implementation handles this without stalling: a read that arrives too early
is registered with \code{read\_index = None}, and \code{\_set\_commit\_index}
assigns it the moment the first entry of the leader's own term (the no-op)
commits:
\end{keyidea}

\begin{lstlisting}[language=Python, caption={raft/node.py: pending reads get their index the instant the term's first entry commits}]
if self.role is Role.LEADER and self._term_commit_index is None and \
        self.log.term_at(index) == self.current_term:
    self._term_commit_index = index
    for r in self._pending_reads:
        if r.read_index is None:
            r.read_index = self.commit_index
\end{lstlisting}

\code{docs/ARCHITECTURE.md} also documents the alternative that was rejected:
\emph{leader leases}, which let a leader serve reads for
\code{election\_timeout - clock\_drift} after a heartbeat quorum, saving the
per-read round trip. The stated reason for refusing it is exact and worth
repeating in an interview: leases buy latency at the price of a safety argument
that depends on \textbf{bounded clock drift between physical machines}. In a
project whose entire thesis is testable correctness, importing an untestable
assumption about hardware clocks would be self-defeating. ReadIndex costs one
heartbeat round and assumes nothing about clocks.

A pleasing efficiency: the heartbeat round is shared. Concurrent reads
piggyback on \code{\_hb\_seq} windows, and \code{\_ack\_seq[p]} is monotone
(\code{max(...)}), so an ack for round 9 automatically satisfies a read waiting
on round 7.

% =====================================================================
\section{The KV state machine and exactly-once writes}

\file{kv/machine.py} is 72 lines and is the deterministic program from the
replicated state machine definition. Four operations:

\begin{center}
\begin{tabular}{lll}
\toprule
\textbf{Op} & \textbf{Effect} & \textbf{Result dict} \\
\midrule
\code{put} & \code{data[key] = value} & \code{\{ok: True\}} \\
\code{delete} & remove key & \code{\{ok: True, existed: bool\}} \\
\code{cas} & set only if current value == expected & \code{\{ok, swapped: bool, value: old\}} \\
\code{get} & never in the log; via ReadIndex & \code{\{value\}} \\
\bottomrule
\end{tabular}
\end{center}

\begin{jargon}{Compare-and-swap (cas)}
``Set \code{key} to \code{new}, but only if it currently holds \code{expected}.''
It is the primitive that makes safe concurrent updates possible without locks:
a client can read a value, compute a new one, and cas it, and the cas fails
harmlessly if anybody else changed the key in between. In this codebase, a cas
with \code{expected=None} means ``only if the key is absent'', and a cas with
\code{value=None} means ``delete if it currently equals expected''.
\end{jargon}

\subsection{The exactly-once problem}

\begin{keyidea}{Why retries silently corrupt data, and why the fix must live in the state machine}
A client sends \code{cas team: raptors $\rightarrow$ dinos}. The leader commits it,
then dies before the reply is sent. The client times out and retries against
the new leader. The new leader has no idea this is a retry: it applies the cas
a second time. Now the value is \code{raptors}... no: the value is already
\code{dinos}, so the second cas \emph{fails}, and the client is told its
operation failed when in fact it succeeded. For a non-idempotent operation like
``increment'', a retry would double-apply outright.

The fix: every write carries \code{(client\_id, seq)}. The state machine keeps,
per client, the last seq it applied and the \emph{result it returned}. A
duplicate returns the cached result without re-executing.

The part that is easy to get wrong: this table must live \textbf{inside the
replicated state}, replicated through the log and included in snapshots. If it
were a leader-local cache, failover would silently downgrade the system from
exactly-once to at-least-once, precisely when it matters most, because the new
leader would have no memory of what the old one had already done.
\end{keyidea}

\begin{lstlisting}[language=Python, caption={kv/machine.py: the entire dedup mechanism, at the top of apply()}]
def apply(self, entry: Entry) -> Result:
    cmd = entry.data
    client, seq = cmd.get("client"), cmd.get("seq")
    if client is not None and seq is not None:
        cached = self.sessions.get(client)
        if cached is not None and seq <= cached[0]:
            return cached[1]          # duplicate: original result, no re-execution
    result = self._execute(cmd)
    if client is not None and seq is not None:
        self.sessions[client] = (seq, result)
    return result
\end{lstlisting}

And the snapshot carries it, which is the point:

\begin{lstlisting}[language=Python, caption={kv/machine.py: sessions ride inside the snapshot, not beside it}]
def snapshot(self) -> str:
    return json.dumps({"data": self.data, "sessions": self.sessions})
\end{lstlisting}

\code{test\_exactly\_once\_across\_leader\_failover} drives exactly the scenario
above: commit a cas, crash the leader before the client hears back, retry the
identical \code{(client, seq)} on the new leader, and assert the retry returns
\code{swapped: True} (the \emph{cached original} outcome, not the correct-but-wrong
answer a fresh evaluation would give) and that the value was applied once.

\begin{honest}{The dedup rule is \code{seq <= cached[0]}, which is stricter than it looks}
Because the check is ``less than or equal to the highest seq I have seen'', a
write that arrives with an \emph{older} seq than one already applied is
silently discarded as a duplicate, returning the newer operation's cached
result. That is correct only if each client's seqs are strictly monotonic. The
README describes this biting the project's own test harness: a helper issued
writes with random seq numbers, they were correctly discarded, and two
unrelated tests failed mysteriously. The lesson was baked into \code{client.py}
(\code{self.\_seq = itertools.count(1)}) and into the simulator's client (which
uses the globally monotonic \code{op.op\_id} as its seq). It is worth being
clear about the shape of the contract: \textbf{the correctness of the dedup
table is a constraint imposed on every client, enforced nowhere in the server}.
A client library with a bug here corrupts its own results and nothing detects
it.
\end{honest}

\begin{honest}{Sessions never expire, and the leak is unbounded in the wrong dimension}
\code{self.sessions[client] = (seq, result)} is never removed. The README lists
this. What the README does not spell out is that the table is not just a memory
leak in RAM: it is part of every snapshot, so it is written to disk, shipped
over the network in full on every InstallSnapshot, and grows with the number of
\emph{distinct client ids ever seen}. \code{RaftClient} generates a fresh random
\code{cli-<uuid>} per instance by default, and the CLI constructs a new client
per command, so \textbf{every single \code{raftkv put} from the command line
permanently adds a row}. The dissertation's fix (a session TTL agreed through
the log) is correctly identified as the prerequisite for calling this
production-ready.
\end{honest}

% =====================================================================
\section{The transports}

\subsection{Real TCP}

\begin{jargon}{Length-prefixed framing}
TCP is a stream of bytes with no message boundaries: send two 100-byte messages
and the receiver may get 200 bytes at once, or 3 then 197. So every message is
prefixed with its length (here 4 bytes, big-endian), and the reader reads
exactly 4 bytes, decodes the length, then reads exactly that many. This
codebase uses \code{readexactly()} for both, which is the correct primitive.
\end{jargon}

Design notes worth defending:

\begin{itemize}
\item \textbf{One port for peer and client traffic}, distinguished by a
      \code{kind} field in the frame (\code{"hello"}, \code{"raft"}, or
      anything else meaning a client request).
\item \textbf{Peer sends are deliberately lossy.} Each peer has a bounded
      \code{asyncio.Queue(maxsize=256)} drained by a dedicated writer task
      that reconnects on failure. If the queue is full,
      \code{except asyncio.QueueFull: pass}, and the message is dropped without
      comment. This looks like a bug and is a design decision: Raft already
      assumes a lossy network and retransmits on the next heartbeat, so
      dropping is strictly better than growing an unbounded queue of stale
      messages that will be obsolete by the time they are sent.
\end{itemize}

\begin{keyidea}{The Python 3.12 shutdown deadlock (a real bug, worth telling)}
Python 3.12 changed \code{asyncio.Server.wait\_closed()} to also wait for
in-flight connection \emph{handlers}, not just the listening socket. Every
client connection in this design sits parked in \code{readexactly}, waiting
forever for a next request that will never come. So \code{wait\_closed()} never
returned: the TCP integration test hung for minutes at approximately zero CPU,
which is the signature of a deadlock rather than a slow loop. The fix is the
\code{self.\_conns} set: track every accepted \code{StreamWriter} and close
them all \emph{before} awaiting \code{wait\_closed()}, so the parked handlers
get an EOF, return, and let the server finish closing.
\end{keyidea}

\subsection{The simulated network}

\code{SimNetwork} implements the same send/deliver contract over the virtual
clock, with four seeded faults. The whole of \code{send} is:

\begin{lstlisting}[language=Python, caption={sim/network.py: every fault, in nine lines}]
def send(self, src: str, dest: str, msg: Message) -> None:
    self.messages_sent += 1
    if self.rng.random() < self.drop_prob:          # 1. loss
        self.messages_dropped += 1
        return
    copies = 2 if self.rng.random() < self.dup_prob else 1   # 2. duplication
    for _ in range(copies):
        delay = self.rng.uniform(self.min_delay, self.max_delay)  # 3. reordering, via delay
        self.clock.call_later(delay, lambda: self._deliver(src, dest, msg))

def _deliver(self, src: str, dest: str, msg: Message) -> None:
    handler = self._handlers.get(dest)
    if handler is None or self.is_partitioned(src, dest):   # 4. partition, at DELIVERY time
        self.messages_dropped += 1
        return
    handler(src, msg)
\end{lstlisting}

\begin{keyidea}{Why partitions are checked at delivery, not at send}
Because that is what a real network does. If a partition forms while a message
is in flight, the message dies in transit; the TCP connection carrying it
resets. Checking at send time would model a network that politely refuses to
accept doomed messages, which no network does, and would miss the entire class
of bugs where a message crosses a boundary that closed behind it. Independent
random delay per copy also means duplicates arrive out of order, for free.
\end{keyidea}

\code{SimEnv} adds two things the real \code{AsyncioEnv} does not need:

\begin{itemize}
\item \textbf{\code{timer\_scale}}, multiplying every scheduled delay, which
      models clock skew. A node with scale 0.5 fires its election timers twice
      as fast as everyone else; the nemesis rolls scales in [0.5, 1.8].
\item \textbf{\code{kill()} and the epoch guard.} A simulated crash must
      invalidate every callback the dead node scheduled, including any already
      sitting on the shared heap. Every callback is wrapped in a closure that
      compares the epoch it was created in against the current one and silently
      does nothing if they differ. This is how ``the process died'' is modelled
      without unwinding the heap.
\end{itemize}

\begin{lstlisting}[language=Python, caption={sim/network.py: the epoch guard, i.e. how a node stops existing}]
def call_later(self, delay: float, fn: Callable[[], None]) -> SimTimer:
    epoch = self._epoch
    def guarded() -> None:
        if epoch == self._epoch:
            fn()
    t = self.clock.call_later(delay * self.timer_scale, guarded)
    ...
def kill(self) -> None:
    """Invalidate all outstanding timers (simulated crash)."""
    self._epoch += 1
    ...
\end{lstlisting}

% =====================================================================
\section{The deterministic simulator}

\subsection{The virtual clock}

\code{SimClock} is 57 lines and is the beating heart of the determinism claim.
One binary heap of events, ordered by \code{(when, seq)}. \code{step()} pops the
earliest, jumps \code{now} to its timestamp, and runs it.

\begin{keyidea}{Why the tie-break on a sequence number is load-bearing}
Two events scheduled for the same virtual instant must still run in a defined
order, or determinism dies. \code{SimTimer.\_\_lt\_\_} compares
\code{(self.when, self.seq)}, and \code{seq} is a monotonically increasing
counter assigned at scheduling time. Without it, Python's heap would break ties
by whatever happened to be in the array, and the same seed could produce two
different runs. The README describes exactly this class of bug: two accidental
uses of ambient ordering, caught by a test that runs a seed twice and compares
full history fingerprints.
\end{keyidea}

And because time is a variable rather than a thing you wait for, ``12 virtual
seconds of a 3-node cluster'' costs roughly 0.1 real seconds. There is no
\code{sleep} anywhere in the simulator. That single fact is what converts
``test the distributed system'' from an overnight job into a function call.

\subsection{One simulation run}

\begin{figure}[H]
\centering
\resizebox{\textwidth}{!}{
\begin{tikzpicture}[node distance=7mm]
\node[box, minimum width=26mm] (seed) {\textbf{seed} (one int)};
\node[dec, right=12mm of seed, minimum width=16mm] (root) {root\\Random};
\draw[flow] (seed) -- (root);

\node[extbox, above right=16mm and 14mm of root, minimum width=30mm] (r1) {network rng\\\scriptsize{}delay, loss, dup};
\node[extbox, right=14mm of root, minimum width=30mm, yshift=6mm] (r2) {cluster rng\\\scriptsize{}$\rightarrow$ per-node env rngs\\\scriptsize{}(election jitter)};
\node[extbox, right=14mm of root, minimum width=30mm, yshift=-8mm] (r3) {nemesis rng\\\scriptsize{}what fault, when, to whom};
\node[extbox, below right=16mm and 14mm of root, minimum width=30mm] (r4) {3 x client rng\\\scriptsize{}which key, op, retry timing};
\draw[dflow] (root) -- (r1); \draw[dflow] (root) -- (r2);
\draw[dflow] (root) -- (r3); \draw[dflow] (root) -- (r4);

\node[box, minimum width=32mm, minimum height=30mm, right=16mm of r2, yshift=-6mm] (heap)
  {\textbf{SimClock}\\one heap ordered\\by (time, seq)\\\scriptsize{}timers, deliveries,\\\scriptsize{}client actions, faults};
\draw[flow] (r1) -- (heap); \draw[flow] (r2) -- (heap);
\draw[flow] (r3) -- (heap); \draw[flow] (r4) -- (heap);

\node[box, minimum width=30mm, right=14mm of heap, yshift=10mm] (nodes) {3 x RaftNode\\\scriptsize{}the real thing};
\node[box, minimum width=30mm, right=14mm of heap, yshift=-10mm] (mon) {Monitor\\\scriptsize{}4 invariants,\\\scriptsize{}checked live};
\draw[flow] (heap) -- (nodes);
\draw[flow] (nodes) -- node[lbl, right]{trace events} (mon);

\node[store, minimum width=28mm, below=12mm of heap] (hist) {HistoryRecorder\\\scriptsize{}every invoke/complete};
\draw[flow] (r4) to[out=280, in=180] (hist);
\node[box, minimum width=30mm, right=14mm of hist] (chk) {Linearizability\\checker};
\draw[flow] (hist) -- node[lbl, above]{at the end} (chk);
\node[box, minimum width=24mm, right=12mm of chk] (out) {PASS / FAIL\\\scriptsize{}with the seed};
\draw[flow] (chk) -- (out);
\draw[flow] (mon) to[out=0, in=60] (out);
\end{tikzpicture}}
\caption{One seeded run. The seed fans out into six independent generators, all
events funnel through one ordered heap, and two independent oracles (the live
invariant Monitor and the post-hoc linearizability checker) judge the result.
Nothing anywhere consults a wall clock or the global \code{random} module.}
\end{figure}

The default \code{SimConfig}: 3 nodes, 12 virtual seconds, 3 concurrent
clients, 3 keys, 5\% message loss, 5\% duplication, 1 to 30 ms delay, a nemesis
action every 1.5 virtual seconds on average, snapshot threshold 40 (chosen low,
so that compaction and InstallSnapshot actually happen within 12 seconds rather
than being theoretically reachable).

Three keys and three clients is a deliberate, well-judged choice: it maximizes
\emph{contention} (many operations racing on the same key) while keeping the
per-key history narrow enough for the checker's exponential search to finish.
A thousand keys would produce more operations and test far less.

\subsection{The nemesis}

\begin{jargon}{Nemesis}
The component whose job is to break things: the adversary. The term comes from
Jepsen, the best-known distributed-systems testing tool.
\end{jargon}

Faults fire at exponentially distributed intervals (\code{rng.expovariate}),
which models a Poisson process, i.e. events that arrive independently at a
constant average rate. The action is chosen by one roll:

\begin{center}
\begin{tabular}{llp{7.4cm}}
\toprule
\textbf{Prob} & \textbf{Action} & \textbf{What it does} \\
\midrule
0.35 & partition & heal everything, then shuffle the nodes and cut at a random point, so any bisection is possible \\
0.25 & crash & pick a live node, destroy it (volatile state gone, \code{MemoryStorage} kept, modelling an fsync'd disk), restart it after 0.3 to 2.0 virtual seconds \\
0.15 & skew & set a random node's \code{timer\_scale} to a value in [0.5, 1.8] \\
0.25 & heal & clear all partitions \\
\bottomrule
\end{tabular}
\end{center}

Two details that make the harness honest rather than merely violent:

\begin{itemize}
\item \code{if len(alive) > 1} before crashing: at least one node always
      survives. A cluster where everything is dead is not a test.
\item Chaos stops at 70\% of the run (\code{chaos\_deadline}), and
      \code{\_stabilize()} heals the network, restarts every node, and resets
      every clock skew. Clients then have a calm window (plus a 3-second drain
      after the nominal duration) to finish their in-flight retries.
\end{itemize}

\begin{keyidea}{Why the stabilize window is not cheating}
Without it, the run would end mid-chaos and most operations would be left
incomplete. Incomplete operations are the \emph{easiest} thing for a
linearizability checker to accept, because it can always declare that they
never happened. A harness that ends in chaos would therefore look clean while
testing very little. Stabilizing forces operations to actually complete with
real observed results, which is what gives the checker something to refute. It
is the opposite of cheating: it is what makes the check bite.
\end{keyidea}

\subsection{The four invariants, checked live}

\code{RaftNode} exposes an optional \code{trace} callback (\code{None} in
production, so the cost is one \code{if} per event). The \code{Monitor} hooks
it and checks four properties continuously, each failure naming the seed:

\begin{enumerate}
\item \textbf{Election safety}: at most one leader per term. A dict
      \code{leaders\_by\_term} is kept for the \emph{entire run}, so a second
      leader for term 5 is caught even if it appears long after the first one
      is gone. Per-sample checking would miss this.
\item \textbf{Leader completeness}: the instant any node wins an election, its
      log is checked against every \code{(index, term)} that any node ever
      applied. If a new leader is missing a committed entry, the run fails
      immediately, at the moment of the crime rather than at the eventual
      symptom.
\item \textbf{State machine safety}: every applied index and entry pair is
      recorded globally, in the monitor's \code{applied} dict. Two nodes
      applying different entries at the same index is an immediate failure.
      This is the direct detection of divergence.
\item \textbf{Log matching}: every 200 events, all pairs of live logs are
      compared: entries with equal index and term must be identical.
\end{enumerate}

\begin{honest}{Log matching is checked shallowly, and the comment admits it}
\code{check\_log\_matching} walks each pair of logs downward from the highest
common index and \code{break}s at the first index where the terms agree,
commenting ``matching here implies matching below (checked lazily)''. That
inference is exactly the Log Matching Property, so it is sound \emph{if the
property holds}, which is what is being tested. The check is therefore
partially circular: a bug that broke the prefix invariant beneath a
coincidentally matching entry would not be caught by this check. It would still
be caught by the state-machine-safety invariant and the linearizability checker,
so the safety net has other layers, but this specific check is weaker than the
name suggests.
\end{honest}

\subsection{The canary: testing the test}
\label{sec:canary}

\begin{keyidea}{A safety net that never fires is indistinguishable from one that cannot fire}
This is the single best idea in the test suite.
\code{test\_broken\_raft\_is\_caught} uses \code{monkeypatch} to replace
\code{RaftNode.\_majority} with \code{max(1, len(config) // 2)}, that is, it
deliberately breaks the quorum rule so split brain becomes possible. It then
asserts that sweeping at most 40 seeds detects an invariant violation, and that
the message names the seed. If the harness could not catch a cluster with the
majority rule removed, every green run it ever produced would be worthless.
\end{keyidea}

\begin{lstlisting}[language=Python, caption={tests/test\_sim\_invariants.py: sabotage the algorithm, demand the harness notices}]
def test_broken_raft_is_caught(monkeypatch) -> None:
    monkeypatch.setattr(RaftNode, "_majority",
                        lambda self: max(1, len(self.config) // 2))
    caught = None
    for seed in range(40):
        try:
            run_simulation(SimConfig(seed=seed))
        except InvariantViolation as exc:
            caught = exc
            break
    assert caught is not None, "weakened quorum was never detected"
    assert "seed=" in str(caught)
\end{lstlisting}

The same instinct produces \code{test\_same\_seed\_same\_run}, which runs a seed
twice and compares a full fingerprint (messages sent, ops completed, and the
complete tuple of every operation's client, key, kind, invoke time, complete
time, and result). Determinism is not assumed; it is asserted, structurally,
every test run.

\subsection{The batch runner}

\file{sim/batch.py} maps seeds across a \code{multiprocessing.Pool}. This is a
legitimate use of parallelism that does not compromise determinism, because
each seed is an independent pure function: process 3 running seed 1337 gets
byte-identical results to a single-threaded run of seed 1337. Parallelism
changes the wall-clock cost and nothing else. Failures print with their seed
immediately (\code{flush=True}) and are re-listed at the end.

% =====================================================================
\section{The linearizability checker}

This is written from scratch and is, algorithmically, the most demanding code
in the repository.

\subsection{What it is given}

A \emph{history}: a list of \code{Op} records, each with an invoke time and
either a completion time plus the observed result, or \code{math.inf} for
completion, meaning the client never heard back.

\begin{keyidea}{Incomplete operations are the interesting ones}
If a client's request times out, the operation might have been applied, or
might not have. \emph{Nobody knows, including the database.} A checker that
ignored them would miss real bugs; a checker that assumed they happened would
report false violations. The only correct treatment is to let the search branch
both ways: an incomplete op may be linearized (with any result accepted, since
none was observed) or left out entirely.
\end{keyidea}

\subsection{The algorithm}

\begin{jargon}{Wing \& Gong search, with Lowe's memoization}
The classical decision procedure. Repeatedly pick some operation that is
\emph{minimal} (no un-linearized operation completed before it was invoked, so
real-time order permits it to go next), apply it to a model of a single
register, and check the model reproduces the result the client actually saw. If
yes, recurse; if the branch dies, backtrack. Lowe's contribution is to memoize
visited states so the same (set of linearized ops, register value) pair is never
explored twice, which turns a factorial blowup into something that usually
finishes.
\end{jargon}

\begin{keyidea}{The two decompositions that make this tractable at all}
\textbf{One: per key.} KV keys are independent registers, and no operation in
this system touches two keys, so a history is linearizable if and only if each
key's sub-history is. \code{check\_history} splits by key and checks each
separately. This turns one giant exponential problem into several small ones,
and the README correctly calls it the largest practical win.

\textbf{Two: the bitmask.} The set of already-linearized operations is
represented as an integer where bit \code{i} means ``op \code{i} is done''.
Set membership, insertion, and the memo key all become single machine
instructions on an integer. This is why the checker's limit is \textbf{63
operations per key}: a Python int is arbitrary precision, but the code
deliberately caps it, and returns \code{UNKNOWN} above that rather than
attempting a search that cannot finish.
\end{keyidea}

\begin{lstlisting}[language=Python, caption={checker/linearizability.py: the search loop, in full}]
seen: set[tuple[int, str | None]] = set()
explored = 0
# Depth-first search over (linearized-mask, register value).
stack: list[tuple[int, str | None]] = [(0, _ABSENT)]
seen.add((0, _ABSENT))
while stack:
    mask, value = stack.pop()
    if mask & complete_mask == complete_mask:
        # Every completed op linearized; leftovers are incomplete ops
        # we are free to declare "never happened".
        return CheckResult(Outcome.OK, key, explored)
    # An op is minimal if no un-linearized op completed before its invoke.
    min_complete = math.inf
    for i, op in enumerate(ops):
        if not mask & (1 << i):
            min_complete = min(min_complete, op.complete)
    for i, op in enumerate(ops):
        bit = 1 << i
        if mask & bit or op.invoke > min_complete:
            continue
        matches, new_value = _step(value, op)
        if not matches:
            continue
        state = (mask | bit, new_value)
        if state not in seen:
            seen.add(state)
            explored += 1
            if explored > max_states:
                return CheckResult(Outcome.UNKNOWN, key, explored,
                                   "state budget exhausted")
            stack.append(state)
return CheckResult(Outcome.VIOLATION, key, explored, ...)
\end{lstlisting}

Reading that loop carefully repays the effort:

\begin{itemize}
\item The success condition is \code{mask \& complete\_mask == complete\_mask},
      not \code{mask == all ones}. It requires every \emph{completed} op to be
      linearized, and lets any leftover incomplete ops be declared to have
      never happened. That is the ``branch both ways'' semantics, implemented
      not as a branch but as a weaker acceptance test, which is neater and
      cheaper.
\item \code{min\_complete} is the earliest completion among
      \emph{un-linearized} ops. An op is a legal next choice if it was invoked
      at or before that time. This is exactly the real-time-order constraint.
\item Falling out of the \code{while} loop means the search space was exhausted
      with no explanation found, and \emph{that} is a violation: not a timeout,
      a proof.
\end{itemize}

\begin{keyidea}{Three outcomes, never two}
\code{OK}, \code{VIOLATION}, and \code{UNKNOWN}. The third is the honest one.
Linearizability checking is NP-complete in general, so any checker must
eventually refuse. This one refuses in two explicit places (more than 63 ops on
a key; more than \code{max\_states} states explored, default 2,000,000) and
returns \code{UNKNOWN}, which is never reported as a pass and never as a fail.
The design property that matters: \textbf{the checker never returns a false
PASS or a false FAIL}, only sometimes a shrug.
\end{keyidea}

\begin{figure}[H]
\centering
\resizebox{0.95\textwidth}{!}{
\begin{tikzpicture}[node distance=9mm]
\tikzset{st/.style={draw=inkblue, fill=softgrey, rounded corners=2pt, font=\scriptsize, inner sep=3pt, align=center}}
\tikzset{dead/.style={st, draw=warnred, fill=warnred!8, text=warnred}}
\tikzset{good/.style={st, draw=okgreen, fill=okgreen!12, text=okgreen}}

\node[st] (r) {mask=000\\value=absent};
\node[st, below left=11mm and 22mm of r] (a) {mask=001\\value="a"\\\scriptsize{}put a};
\node[dead, below=11mm of r] (b) {mask=010\\get$\rightarrow$"a"\\but value is absent:\\\textbf{prune}};
\node[st, below right=11mm and 22mm of r] (c) {mask=100\\value="b"\\\scriptsize{}put b};

\draw[flow] (r) -- node[lbl, above, sloped]{op0: put "a"} (a);
\draw[flow, draw=warnred] (r) -- node[lbl, right]{op1: get $\rightarrow$ "a"} (b);
\draw[flow] (r) -- node[lbl, above, sloped]{op2: put "b"} (c);

\node[good, below=11mm of a] (a2) {mask=011\\value="a"};
\node[dead, below right=11mm and 2mm of a] (a3) {mask=101\\value="b"\\\scriptsize{}memoized: seen};
\draw[flow, draw=okgreen] (a) -- node[lbl, left]{op1: get $\rightarrow$ "a"\\model agrees} (a2);
\draw[flow, draw=warnred] (a) -- node[lbl, right, pos=0.7]{op2} (a3);
\node[good, below=9mm of a2] (a4) {mask=111\\\textbf{all completed ops}\\\textbf{linearized: OK}};
\draw[flow, draw=okgreen] (a2) -- (a4);

\node[lbl, anchor=west, align=left, text width=5.4cm] at (4.6,-3.4)
  {Each node is a memo key: (mask, value).\\
   Pruned when the model contradicts the\\
   observed result, or when the state was\\
   already visited by another path.\\
   Exhausting the tree with no OK leaf\\
   is a \textbf{proof} of violation, not a timeout.};
\end{tikzpicture}}
\caption{The search, on a three-operation history. Branches die either because
the model cannot reproduce what the client saw, or because Lowe memoization has
already visited that exact state. Both prunings are what keep the exponential
in check.}
\end{figure}

\subsection{The checker's own tests}

\code{test\_linearizability.py} is 11 tests, and the important ones are the
ones that demand \emph{rejection}:

\begin{description}[leftmargin=0pt, itemsep=4pt]
\item[\code{test\_stale\_read\_is\_violation}]\hfill\\
      a put completed strictly before a get began, yet the get missed it.
\item[\code{test\_value\_from\_nowhere\_is\_violation}]\hfill\\
      a get returned a value nobody ever wrote.
\item[\code{test\_lost\_update\_is\_violation}]\hfill\\
      two non-overlapping cas from the same expected value both succeeded.
\item[\code{test\_known\_bad\_interleaving\_from\_paper\_shape}]\hfill\\
      a read after a completed cas returned the pre-cas value.
\end{description}

\begin{keyidea}{The budget test that passed for the wrong reason: the best story in the repo}
The README tells this on itself, and it deserves the space. The original test
for ``does the budget cause UNKNOWN'' used a history of 40 incomplete
concurrent puts. It passed. But it passed for a reason that had nothing to do
with the budget: a history of purely \emph{incomplete} operations is trivially
linearizable by declaring that none of them ever happened, so the search
succeeded on its very first move and never explored anything. The test asserted
the right outcome while exercising none of the mechanism.

The fixed test buries a real violation (\code{get $\rightarrow$ "phantom"})
under 20 \emph{completed}, fully concurrent puts, which forces the exhaustive
search that the budget exists to bound. This is a precise instance of the
general lesson that runs through the whole project: a test that passes tells
you nothing until you know \emph{why} it passes.
\end{keyidea}

\begin{honest}{The checker validates get and cas results, but not put or delete}
Look at \code{\_step}. For \code{"get"} and \code{"cas"} it compares the model's
answer against \code{op.result}. For \code{"put"} it returns \code{(True,
new\_value)} unconditionally, and for \code{"delete"} it returns \code{(True,
\_ABSENT)} unconditionally. So the \code{existed} flag that \code{delete}
returns is \textbf{never checked against the model}, and a database that
reported \code{existed: False} for a delete of a present key would sail through.
This is a defensible narrowing (put and delete have no interesting return
value, and \code{existed} is arguably not part of the register abstraction) but
it is unstated, and the honest framing is that the checker verifies the
register semantics of get and cas, and treats put and delete as pure state
transitions.
\end{honest}

\begin{honest}{The chaos runs never exercise delete at all}
\code{SimClient.\_next\_op} rolls: below 0.4, get; below 0.75, put; otherwise
cas. There is no branch that issues a \code{delete}. So the headline sweep,
however many seeds it runs, tests \code{delete} exactly zero times end to end.
\code{delete} is covered by \code{test\_basic\_ops} (a unit test on
\code{KVMachine} alone) and by one line of the TCP integration test, and that
is all. It is a small gap, and easy to close (one branch in the client), but a
claim of the form ``564,170 operations verified'' should be read as ``get, put
and cas operations''.
\end{honest}

% =====================================================================
\section{Verification: what was actually run for this document}
\label{sec:verified}

The brief for this document required that the project's headline claims be
\textbf{executed rather than repeated}. They were, on this machine, from a
clean virtual environment (\code{python3.12 -m venv .venv} and
\code{pip install -e ".[dev]"}), on Linux with Python 3.12.3, 8 CPU cores.

\subsection{The test suite}

\begin{lstlisting}[caption={Observed output: the full suite, including the slow 200-seed sweep}]
$ .venv/bin/python -m pytest
============================= test session starts ==============================
platform linux -- Python 3.12.3, pytest-9.1.1, pluggy-1.6.0
rootdir: /home/rolex/Salman Adnan/Programming/Portfolio/raft-kv
configfile: pyproject.toml
testpaths: tests
collected 57 items

tests/test_election.py .....                                             [  8%]
tests/test_kv_sessions.py ....                                           [ 15%]
tests/test_linearizability.py ...........                                [ 35%]
tests/test_membership.py .....                                           [ 43%]
tests/test_persistence.py .......                                        [ 56%]
tests/test_readindex.py ....                                             [ 63%]
tests/test_replication.py .........                                      [ 78%]
tests/test_sim_invariants.py ......                                      [ 89%]
tests/test_snapshot.py ...                                               [ 94%]
tests/test_tcp_integration.py ...                                        [100%]

============================= 57 passed in 51.41s ==============================
\end{lstlisting}

\textbf{Verified.} 57 tests, all passing, including the three real-TCP
integration tests that bind localhost sockets and the 200-seed slow sweep. The
README's claim of ``57 tests, all passing'' is exact.

\subsection{The headline sweep}

\begin{lstlisting}[caption={Observed output: the 5000-seed sweep, run for this document}]
$ .venv/bin/raftkv sim --seeds 5000 --workers 8
... 250/5000 seeds (13s)
... 500/5000 seeds (27s)
... (elided: one progress line per 250 seeds)
... 4750/5000 seeds (243s)
... 5000/5000 seeds (257s)
5000 seeds in 257.0s (19.5 seeds/s), 564170 client ops total, 0 failure(s)
\end{lstlisting}

\begin{keyidea}{The op count reproduced exactly, and that is the real result}
The README reports \code{564170 client ops total} from a run with
\code{--workers 10}. This machine ran it with \code{--workers 8}. The op count
came back \textbf{identical to the digit}. That is not a coincidence and it is
not a weak check: it is the determinism claim being confirmed end to end. If any
part of the system had consulted a wall clock, a hash seed, a set iteration
order, or the global \code{random} module, running the same 5000 seeds on a
different day with a different worker count would have produced a different
number. It did not. The timing legitimately differs (fewer workers, different
machine, different load), and timing is the one thing that is allowed to
differ, because it is the only thing outside the seed.
\end{keyidea}

\begin{center}
\begin{tabular}{llll}
\toprule
\textbf{Claim (README)} & \textbf{Observed here} & \textbf{Verdict} \\
\midrule
57 tests, all passing & 57 passed in 51.41s & confirmed \\
5000 seeds & 5000 seeds & confirmed \\
564170 client ops total & see above & confirmed, exactly \\
0 failures & 0 failure(s) & confirmed \\
84.7s / 59.0 seeds/s at \code{--workers 10} & 257.0s / 19.5 seeds/s at \code{--workers 8} & differs, see below \\
\bottomrule
\end{tabular}
\end{center}

\begin{honest}{The one number that did not reproduce is the timing, and the gap is bigger than the worker count explains}
The README reports 84.7 seconds (59.0 seeds/s) with \code{--workers 10}. This
run took 257.0 seconds (19.5 seeds/s) with \code{--workers 8}. Dropping from 10
workers to 8 accounts for at most a factor of 1.25; the observed gap is a factor
of about 3.

This is not a discrepancy in the code, and it is not evidence against anything:
throughput is the one property of the sweep that is deliberately \emph{outside}
the seed, and it depends on the core count, the CPU, and whatever else the
machine was doing. \textbf{This particular run shared the machine with other
substantial workloads}, which is the most likely explanation, but that is an
\textbf{inference}: it was not isolated or measured, and the README's original
run was not repeated under matched conditions. The honest position is that the
timing figure in the README is a machine-and-day observation rather than a
property of the project, and that the two figures that \emph{are} properties of
the project, the op count and the failure count, both reproduced exactly.
\end{honest}

\subsection{What the numbers do and do not mean}

\begin{honest}{Read the headline number carefully: it is impressive, and it is narrower than it sounds}
``5,000 randomly generated disasters and 564,170 reads/writes with 0
contradictions'' is true as stated and was reproduced here. Being precise about
its scope is what makes it defensible under questioning:

\begin{itemize}
\item \textbf{5000 seeds is not 5000 independent disaster shapes.} Every seed
      runs the same fixed recipe: 3 nodes, 12 virtual seconds, the same fault
      probabilities. The seed varies the timing, the targets and the ordering,
      not the shape of the experiment. It is a deep search of one scenario, not
      a broad search across scenarios.
\item \textbf{The 564,170 ops are get, put and cas.} No \code{delete} is ever
      issued by the simulated client (see the honest box above).
\item \textbf{The disks are perfect.} The simulator uses
      \code{MemoryStorage}, which models an ideal fsync'd disk that never tears
      a write and never reorders. The excellent torn-tail recovery code in
      \code{FileStorage} is therefore \emph{never once exercised by any of the
      5000 seeds}; it is covered only by targeted unit tests. The README lists
      this under future work, and it is the single largest gap between what the
      simulator tests and what production runs.
\item \textbf{Membership changes are never exercised by the sweep either.} The
      nemesis crashes, partitions and skews, but never adds or removes a server.
\item \textbf{``0 contradictions'' means: no invariant violation fired, and no
      history was refuted.} It does not mean every history was proven
      linearizable, because a history that exceeded the checker's budget would
      return \code{UNKNOWN}, and \code{SimResult.ok} treats \code{UNKNOWN} as
      acceptable (\code{all(r.outcome is not Outcome.VIOLATION ...)}). No
      \code{UNKNOWN} count is reported by the batch runner, so the sweep cannot
      distinguish ``5000 proofs'' from ``4990 proofs and 10 shrugs''. Given the
      configured width (3 keys, 3 clients, bounded concurrency) the histories
      almost certainly fit well inside the budget, and the README says so, but
      the runner does not \emph{show} that, and it easily could.
\end{itemize}
\end{honest}

That last point is the one worth volunteering in an interview before anyone
finds it, because volunteering it demonstrates exactly the discipline the whole
project is arguing for.

% =====================================================================
\section{The outer layers: server, client, CLI}

\subsection{RaftServer}

\file{server.py} glues \code{FileStorage} + \code{KVMachine} + \code{RaftNode} +
\code{TcpTransport} into one node and translates client frames into Raft
operations. The interesting mechanism is how a client request waits for its own
entry to commit: \code{\_register(index, term)} creates an
\code{asyncio.Future}, stores it under the key \code{(index, term)}, and awaits
it with a 5-second timeout. The apply listener resolves it.

\begin{lstlisting}[language=Python, caption={server.py: resolve the waiter for this entry, and fail any waiter whose slot was stolen}]
def _on_apply(self, index: int, term: int, entry: Entry, result: Any) -> None:
    fut = self._pending.pop((index, term), None)
    if fut is not None and not fut.done():
        fut.set_result(result)
    for key in [k for k in self._pending if k[0] == index and k[1] != term]:
        stale = self._pending.pop(key)
        if not stale.done():
            stale.set_result(None)
\end{lstlisting}

Keying by \code{(index, term)} rather than \code{index} alone is the whole
point. If this node's leadership was interrupted and a different leader put a
\emph{different} entry at index 42, then the waiter for \code{(42, term 5)} must
not be resolved by the arrival of \code{(42, term 6)}: its entry was silently
overwritten and its client must be told the write did not happen. That is what
the second half of the function does. \code{\_on\_step\_down} does the same,
wholesale, on losing leadership. \code{SimCluster} carries the same logic in
\code{\_make\_apply\_listener}, deliberately mirrored so the simulator's client
sees the same failure semantics as a real one.

\begin{honest}{One asymmetry between the two mirrored implementations}
\code{server.py::\_on\_apply} guards every resolution with
\code{if not fut.done()}. \code{harness.py::\_make\_apply\_listener} calls
\code{host.pending.pop((i, t))(None)} with no equivalent guard, because its
callbacks are plain functions rather than futures and are made idempotent by a
\code{responded} flag inside \code{SimClient.\_attempt} instead. The two are
each correct in context, but the invariant ``a client callback fires at most
once'' is enforced in a different place in each, which is exactly the sort of
divergence that rots.
\end{honest}

\subsection{RaftClient}

Retries with leader hints. The loop is worth one note: on a \code{not\_leader}
reply, it adopts the hint, but \code{self.\_leader\_hint = hint if hint !=
target else None} guards against a node that names \emph{itself} as leader while
rejecting the request, which would otherwise spin. Writes carry
\code{seq=next(self.\_seq)} from an \code{itertools.count(1)}, which is the
monotonic-counter lesson from the dedup incident encoded permanently into the
library.

\begin{honest}{A connection per request, and a session row per CLI invocation}
\code{\_call\_node} opens a fresh TCP connection, sends one frame, reads one
reply, and closes. The README lists persistent connections with request
multiplexing as unfinished work, and notes the frame protocol was designed for
it (frames carry an \code{id} field which is currently only used to assert the
reply matches). Compounding it: the CLI builds a new \code{RaftClient} per
command, and \code{RaftClient} defaults to a fresh random \code{cli-<uuid>}, so
every \code{raftkv put} both opens a new connection and permanently grows the
replicated session table by one row. For a demo this is invisible. It is worth
saying out loud that the CLI is a demonstration surface, not a client library.
\end{honest}

\subsection{The CLI}

\code{raftkv cluster start -n 3} spawns one real OS process per node with
\code{start\_new\_session=True}, records pids in \code{cluster.json}, and then
polls until a leader exists so that the command does not return a cluster that
is not yet usable. This matters for the demo's credibility: because nodes are
real processes, \code{demo.sh}'s \code{kill -9} on the leader pid is a genuine
uncatchable process kill, not a polite in-process shutdown. \code{SIGKILL}
cannot be trapped, so no cleanup handler runs, no buffers are flushed, and
recovery must work purely from what \code{fsync} had already put on disk.

% =====================================================================
\section{Consolidated honest assessment}

Everything below is drawn from reading the code, not from the README, though
several items appear in the README's own ``honest limits'' section, which is
itself unusually candid.

\subsection{What is genuinely strong}

\begin{itemize}
\item \textbf{The sans-IO boundary is real, not aspirational.} \file{raft/} imports
      \code{enum} and \code{typing} and its own siblings, and nothing else. No
      \code{asyncio}, no \code{time}, no \code{socket}, no global
      \code{random}. The architecture's central claim survives a grep.
\item \textbf{The subtle rules are all present and correctly implemented}: vote
      persistence before reply, the up-to-date check as a tuple comparison, the
      section 5.4.2 own-term commit restriction with the no-op, the conflict
      hint that can only move \code{next\_index} downward, the directory fsync
      after rename, config effective on append, the leader that defers its own
      removal until commit.
\item \textbf{The tests test the tests.} The weakened-quorum canary, the
      known-bad histories, and the determinism fingerprint are three
      independent instances of the same discipline, and it is rare.
\item \textbf{The honesty is load-bearing.} The \code{UNKNOWN} outcome, the
      documented rejection of leader leases with reasons, and the README's own
      list of limits are the work of someone optimizing for being right rather
      than for looking right.
\end{itemize}

\subsection{Weaknesses, ranked by what would hurt most in an interview}

\begin{enumerate}
\item \textbf{The simulator's disk is perfect, so the crash-recovery code is
      the least-tested important code in the project.} \code{MemoryStorage}
      never tears a write, never reorders an fsync, never fails. The CRC
      scan-and-truncate logic is genuinely good and is exercised only by
      hand-built unit tests that write corrupt bytes directly. The project's
      thesis is ``search the schedule space for bugs''; storage faults are
      simply outside the searched space. The README names this. It is the
      answer to ``what would you do next'', and it should be volunteered.
\item \textbf{\code{UNKNOWN} is invisible.} \code{SimResult.ok} accepts it and
      \code{batch.py} never counts it, so the sweep's headline cannot
      distinguish a proof from a shrug. The fix is a few lines and would make
      the headline number strictly stronger.
\item \textbf{Replication is quadratic in the worst case.} Every heartbeat to a
      lagging follower re-sends \code{tuple(self.log.slice\_from(ni))}, the
      \emph{entire} remaining suffix, because \code{next\_index} only advances
      on a reply. With a 30 ms heartbeat and a follower 10,000 entries behind,
      that is the whole 10,000-entry suffix serialized to JSON every 30 ms
      until the first reply lands. It is correct, because AppendEntries is
      idempotent, and the README calls it naive and wasteful, which is fair.
      There is no batch-size cap.
\item \textbf{\code{\_load\_snapshot()} re-reads and re-parses the entire
      persistent state, log included, on every InstallSnapshot send.}
      See the honest box in the snapshot section.
\item \textbf{No pre-vote.} A node returning from a partition has inflated its
      term by campaigning fruitlessly, so on rejoining it forces a real leader
      to step down and the cluster to hold a pointless election. Raft's
      pre-vote extension fixes this; the README lists it, and correctly notes
      that the simulator would verify the fix in minutes.
\item \textbf{Sessions never expire}, and the table is inside the replicated
      state and every snapshot.
\item \textbf{Scale.} Python, JSON on the wire, one fsync per append batch, one
      connection per client request. This is a correctness project. Its own
      README says so in its first honest limit, which is the right call: an
      inflated throughput claim would be the easiest thing in the world to
      puncture.
\end{enumerate}

\subsection{Small things a careful reader will find}

\begin{itemize}
\item The \code{next\_index} initialization disagrees with itself across
      \code{\_become\_leader}, \code{\_change\_config}, and
      \code{\_refresh\_config} (off by one in the middle one; self-correcting).
\item \code{\_config\_from}'s index filter is effectively vacuous at both call
      sites.
\item \code{Op.args} is typed \code{dict[str, Any]} and \code{Op.result} is
      \code{Any}. The checker's \code{\_step} indexes the \code{value} and
      \code{expected} keys of \code{op.args} with no validation, so a
      malformed history raises \code{KeyError} rather than reporting a bad
      input.
\item \code{RaftLog.compact\_to} silently returns when
      \code{index <= snapshot\_index}, which is the right behaviour for an
      idempotent call but hides a genuine ordering error if one ever occurs.
\item \code{cmd\_cluster\_start} returns exit code 1 with a warning if no leader
      appears within 15 seconds, but leaves the spawned processes running. A
      failed start is not cleaned up.
\item \code{MAX\_FRAME} is 64 MB and \code{read\_frame} returns \code{None}
      (silently closing the connection) rather than logging when a frame
      exceeds it. Given uncapped, unchunked snapshots, this is the exact place
      a too-large snapshot would fail, and it would fail quietly.
\end{itemize}

% =====================================================================
\section{What to take away}

\begin{keyidea}{The five sentences that summarize the project}
\textbf{One.} Consensus reduces to keeping an ordered log identical across
machines, and Raft keeps it identical by electing one leader per term who alone
appends, and by requiring a majority for every decision, so that the overlap
between any two majorities makes contradiction impossible.

\textbf{Two.} The single architectural decision everything else rests on is
sans-IO: \code{RaftNode} touches the world only through the \code{Env} and
\code{Storage} protocols, so the identical consensus code runs over real TCP and
inside a simulator, and every simulator bug is by construction a production bug.

\textbf{Three.} Because every source of nondeterminism is funnelled through one
seeded generator and one ordered event heap, a 12-second cluster lifetime under
sustained chaos is a pure function of one integer, which makes thousands of
disaster schedules searchable in seconds and makes any failure reproducible
forever.

\textbf{Four.} Correctness is judged from the outside, by a from-scratch
linearizability checker that either proves a history explainable by a single
honest register, refutes it, or honestly says \code{UNKNOWN}, never guessing.

\textbf{Five.} The safety nets are themselves tested (weaken the quorum rule and
the harness must notice; feed the checker known-bad histories and it must
reject them), because a net that never fires is indistinguishable from one that
cannot.
\end{keyidea}

\subsection{The questions this project invites, and the honest answers}

\begin{center}
\begin{tabular}{p{5.6cm}p{8.6cm}}
\toprule
\textbf{Question} & \textbf{Honest answer} \\
\midrule
Why not use an existing Raft library? & The point was to understand where the subtlety lives. Section 5.4.2 and the ReadIndex term-commit precondition are two places a plausible shortcut is a real bug, and the chaos harness found the second one during development. \\
Is 5000 seeds meaningful? & It is a deep search of one fixed scenario (3 nodes, 12s, fixed fault rates), not a broad one. It reproduced to the digit on a different worker count, which is the determinism claim confirmed. \\
What is not tested? & Storage faults (the simulator's disk is perfect), \code{delete} end to end, membership changes under chaos, and anything above single-digit node counts. \\
What would you do first? & Inject torn writes and fsync reordering into simulated storage, so the 5000 seeds cover the recovery paths that are currently only unit-tested. Then report \code{UNKNOWN} counts in the sweep. \\
Why ReadIndex and not leases? & Leases trade a round trip for an assumption about bounded clock drift between machines. In a project whose thesis is testable correctness, an untestable hardware assumption is a bad trade. \\
\bottomrule
\end{tabular}
\end{center}

\vfill
\begin{center}
\small\itshape
Every claim in section~\ref{sec:verified} was produced by executing the code on
this machine, not by quoting the README. Everything labelled inferred is
labelled inferred. Where the code and the documentation disagreed, the code was
taken as the source of truth.
\end{center}

\end{document}
